% !TeX TXS-program:pdflatex = pdflatex -synctex=1 -recorder -interaction=nonstopmode --shell-escape %.tex
\documentclass[12pt,a4paper]{exam}%Documentation de la classe exam : https://mirrors.ircam.fr/pub/CTAN/macros/latex/contrib/exam/examdoc.pdf
%\usepackage[abspath]{currfile}
\usepackage[utf8]{inputenc}
\usepackage[T1]{fontenc}
\usepackage[french]{babel} 
%\usepackage{amsmath,amssymb}
\usepackage[explicit]{titlesec}
\usepackage{graphicx}
%\usepackage{booktabs}
\usepackage{enumitem}
\usepackage{multirow}
\usepackage{lastpage}
\usepackage{macrolulu2}
\usepackage{tikz}
\usetikzlibrary{calc,shadows}
\usepackage{hyperref}

%\newcommand{\cheminexo}{}
%\newcommand{\cheminexo}{"Z:/Travail/Cours finis/Python/0001 cours perso/2024 2025/L3/Exercices/"}
%\newcommand{\chemincode}{"Z:/Travail/Cours finis/Python/0001 cours perso/2024 2025/L3/Exercices/code/"}
\newcommand{\chemincode}{code/}
%\newcommand{\chemincode}{\currfileabsdir /code/}

\renewcommand{\questionshook}{%
    %\setlength{\leftmargin}{0pt}%
    \setlength{\labelwidth}{-\labelsep}%
    \setlength{\itemsep}{0.5cm}} 
    
    
%En-tête et pied de page
\pagestyle{headandfoot}
\header{\includegraphics[scale=0.035]{python}}{\small\textbf{Banque de questions}}{\includegraphics[scale=0.035]{python}}

%\headrule
\footrule
\footer{}{\scriptsize Page \thepage{} sur \pageref{LastPage}}{}

\titleformat{\section}[display]
  {\normalfont}
  {}
  {0pt}
  {%
    \begin{tikzpicture}
    	\node[minimum width=\textwidth,anchor=west] (B) at (0,0) {};
    	\node[font=\bfseries\large,
    			rounded corners,
    			draw=orange,
    			fill=orange!10,
    			inner sep=8pt,
    			line width=2pt,
    			anchor=west] (A) at (0,0) {\thesection \quad #1};
    	\draw[orange] (A.east) -- (B.east);
    	\foreach \x in {60,45,30,15}
    		{\fill[orange!\x] ([xshift=- \x pt]B.east) circle [radius=10pt];}
    \end{tikzpicture}%
  }
\titlespacing{\section}{0pt}{0pt}{0.5cm}



%\qformat{\textbf{Exercice \thequestion~:}\hfill}
\qformat{\begin{tikzpicture}
    	\node[minimum width=\textwidth,anchor=west] (B) at (0,0) {};
    	\node[font=\bfseries,
    			rounded corners,
    			draw=orange,
    			fill=orange!10,
    			inner sep=6pt,
    			line width=2pt,
    			anchor=west] (A) at (0,0) {Question \thequestion \quad};
    	\draw[orange] (A.east) -- (B.east);
    	\foreach \x in {60,45,30,15}
    		{\fill[orange!\x, drop shadow] ([xshift=- \x pt]B.east) circle [radius=5pt];}
    \end{tikzpicture}
        \vrule depth 1em width 0pt}
%\extrawidth{1cm}
\renewcommand{\solutiontitle}{\noindent\textbf{Code python :}\par\noindent}
%\newenvironment{\reqresult}[2][1]
\newcommand{\img}[2][1]
{
\begin{figure}[H]
	\begin{center}
		\includegraphics[scale=\#1]{\chemincode \#2}
	\end{center}
\end{figure}
}
%baseline={([yshift=-1em] current bounding box.north)}
%baseline=(label.base)
\DeclareRobustCommand{\circled}[1]{%
	\tikz[baseline={([yshift=-1em] current bounding box.north)}]{\node[scale=0.4,circle,
		black, circular drop shadow={shadow xshift=0.2ex, yshift=0.1ex},font=\bfseries,fill= orange!30, draw=orange,inner sep=2pt,minimum size=2.5em] (label) {#1};}}

\setlist[enumerate,1]{
	label =\circled{\Large\arabic*},
	leftmargin=2cm, rightmargin=1cm,
	itemsep=5pt
}
\newcommand*{\MyBall}{\tikz \draw [baseline, ball color=orange!70, draw=orange!20!white] circle (2pt);}
%
\setlist[itemize,1]{label ={\MyBall},leftmargin=2cm, rightmargin=1cm}

%\definecolor{monjaune}{RGB}{255,196,12}
%\definecolor{monjaune}{named}{orange!10}
\colorlet{monjaune}{orange!10}
\definecolor{monmauve}{RGB}{87,23,69}
\definecolor{monnoir}{RGB}{51,51,51}

\newlength{\xE}
\setlength{\xE}{2cm}
\newlength{\xA}
\setlength{\xA}{4cm}
\newlength{\xC}
\setlength{\xC}{13cm}
\newlength{\yE}
\setlength{\yE}{18cm}
\newlength{\yA}
\setlength{\yA}{9cm}
\newlength{\yC}
\setlength{\yC}{6cm}
\newlength{\carreE}
\setlength{\carreE}{5cm}
\newlength{\carreA}
\setlength{\carreA}{12cm}
\newlength{\carreC}
\setlength{\carreC}{6cm}


\printanswers					%Commenter cette ligne pour cacher les solutions.

\newcommand{\nomfichier}{pascorrige.py}

\begin{document}
\thispagestyle{empty}
		\begin{tikzpicture}[
			remember picture,overlay,
			inner sep=0pt,
			outer sep=0pt,
			]
		\coordinate (A) at ([xshift=\xA,yshift=-\yA]current page.north west);
		\coordinate (B) at ([xshift=\xA+\carreA,yshift=-(\yA+\carreA)]current page.north west);
		\coordinate (AA) at ([xshift=\carreA-3.5cm,yshift=0]A);
		\coordinate (BB) at ([xshift=0,yshift=\carreA-3.5cm]B);
		\coordinate (AAA) at ([xshift=0,yshift=-\carreA+3.5cm]A);
		\coordinate (BBB) at ([xshift=-\carreA+3.5cm,yshift=0]B);
		\fill[monjaune,draw=orange,line width= 2pt] (A) -- (AA) -- (BB) -- (B) -- (BBB) -- (AAA) -- cycle;
		
		\coordinate (E) at ([xshift=\xE,yshift=-\yE]current page.north west);
		\coordinate (F) at ([xshift=\xE+\carreE,yshift=-(\yE+\carreE)]current page.north west);
		\coordinate (EE) at ([xshift=\carreE-3cm,yshift=0]E);
		\coordinate (FF) at ([xshift=0,yshift=\carreE-3cm]F);
		\coordinate (EEE) at ([xshift=0,yshift=-\carreE]E);
		\fill[monnoir] (E) -- (EE) -- (FF) -- (F) -- (EEE)  -- cycle;
		
		\coordinate (C) at ([xshift=\xC,yshift=-\yC]current page.north west);
		\coordinate (D) at ([xshift=\xC+\carreC,yshift=-(\yC+\carreC)]current page.north west);
		\coordinate (CC) at ([xshift=0,yshift=-\carreC+3cm]C);
		\coordinate (DD) at ([xshift=-\carreC+3cm,yshift=0]D);
		\coordinate (CCC) at ([xshift=\carreC,yshift=0]C);
		\fill[monnoir] (C) -- (CC) -- (DD) -- (D) -- (CCC)  -- cycle;
		
		%*****bas
		\coordinate (G) at ([yshift=2.5cm]current page.south west);
		\coordinate (H) at ([xshift=4cm,yshift=2.5cm]current page.south);
		\coordinate (I) at ([xshift=3cm,yshift=0.5cm]current page.south);
		\coordinate (J) at ([yshift=0.5cm]current page.south west);
		\coordinate (K) at ([xshift=5cm,yshift=2.5cm]current page.south);
		\coordinate (L) at ([yshift=2.5cm]current page.south east);
		\coordinate (M) at ([yshift=0.5cm]current page.south east);
		\coordinate (N) at ([xshift=4cm,yshift=0.5cm]current page.south);
		\fill[monnoir] (G) -- (H) -- (I) -- (J) -- cycle;
		\fill[orange] (K) -- (L) -- (M) -- (N) -- cycle;		
										
		%*****haut
		\coordinate (GG) at ([yshift=-1cm]current page.north west);
		\coordinate (HH) at ([xshift=-3cm,yshift=-1cm]current page.north);
		\coordinate (II) at ([xshift=-4cm,yshift=-0.5cm]current page.north);
		\coordinate (JJ) at ([yshift=-0.5cm]current page.north west);
		\coordinate (KK) at ([xshift=-3cm,yshift=-1cm]current page.north);
		\coordinate (LL) at ([yshift=-1cm]current page.north east);
		\coordinate (MM) at ([yshift=-0.5cm]current page.north east);
		\coordinate (NN) at ([xshift=-4cm,yshift=-0.5cm]current page.north);
		\fill[monnoir] (GG) -- (HH) -- (II) -- (JJ) -- cycle;
		\fill[orange] (KK) -- (LL) -- (MM) -- (NN) -- cycle;				

		\node [scale=0.7,
				align=center,
				font=\sffamily\bfseries\huge,
				text width=10cm
				] at ([xshift=6cm,yshift=-6cm]A) {
			\includegraphics[scale=0.15]{python}\par
			\vspace{1cm}Python \par
			Banque de questions
			};		
		
		\node [scale=0.6,
				align=center,
				font=\sffamily\bfseries\huge,
				text width=10cm,
				orange
				] at ([xshift=6.5cm,yshift=-1cm]G) {
			Université de TOURS
			};							
		\node [scale=0.5,
				align=center,
				font=\sffamily\bfseries\huge,
				text width=5cm,
				orange
				] at ([xshift=3.5cm,yshift=-2.3cm]C) {
			Frédéric \par LURET
			};	
		\node [scale=0.5,
				align=center,
				font=\sffamily\bfseries,
				text width=10cm,
				align= left,
				anchor = west,
				orange
				] at ([xshift=0.2cm,yshift=0.3cm]EEE) {
			Version du \today
			};					
	\end{tikzpicture}
\clearpage
%\section{Titre}
%\shadedsolutions
\unframedsolutions
%\SolutionEmphasis{\itshape\small}
\tableofcontents


%*********************************************
\section{Site 8 Comprehension Dict}
\href{https://www.tutorjoes.in/python_programming_tutorial/dictionary_comprehensions_exercises_in_python}{Dictionnary de comprehension}\par

\begin{questions}
	
	\renewcommand{\chemincode}{../../code/}
        \question
        DEPARTDICT
Créer un dictionnaire de carrés de nombres de 1 à 10

Exemple de sortie

\{1 : 1, 2 : 4, 3 : 9, 4 : 16, 5 : 25, 6 : 36, 7 : 49, 8 : 64, 9 : 81, 10 : 100\}
        \par
        \begin{solution}
            \renewcommand{\nomfichier}{q660.py}
            \pythonfile{\chemincode \nomfichier}[][\nomfichier]
            Ce code Python crée un dictionnaire nommé carrés où les clés sont des nombres de 1 à 10, et les valeurs sont les carrés de ces nombres. Voici comment fonctionne le code :

    squares = \{x : x ** 2 for x in range(1, 11)\} : Cette ligne utilise une compréhension de dictionnaire pour créer le dictionnaire squares. Elle parcourt chaque nombre x dans l'intervalle de 1 à 10 et affecte une paire clé-valeur au dictionnaire. La clé est le nombre lui-même (x) et la valeur est le carré de ce nombre (x ** 2).
        for x in range(1, 11) : Cette partie du code parcourt chaque nombre de 1 à 10.
    print(squares) : Cette ligne imprime le dictionnaire des carrés sur la console.
        \end{solution}
        

        \question
        Créer un dictionnaire dont les clés sont les nombres pairs et les valeurs leurs carrés.

Exemple de résultat

\{0 : 0, 2 : 4, 4 : 16, 6 : 36, 8 : 64, 10 : 100, 12 : 144, 14 : 196, 16 : 256, 18 : 324, 20 : 400\}
        \par
        \begin{solution}
            \renewcommand{\nomfichier}{q661.py}
            \pythonfile{\chemincode \nomfichier}[][\nomfichier]
            Ce code Python crée un dictionnaire nommé evens\_squared dont les clés sont des nombres pairs compris entre 0 et 20, et dont les valeurs sont les carrés de ces nombres pairs. Voici comment fonctionne ce code :

    evens\_squared = \{x : x ** 2 for x in range(0, 21) if x \% 2 == 0\} : Cette ligne utilise une compréhension de dictionnaire pour créer le dictionnaire evens\_squared. Elle parcourt chaque nombre pair x dans l'intervalle de 0 à 20 et affecte une paire clé-valeur au dictionnaire. La clé est le nombre pair lui-même (x) et la valeur est le carré de ce nombre pair (x ** 2).
        for x in range(0, 21) : Cette partie du code parcourt chaque nombre de 0 à 20.
        if x \% 2 == 0 : elle vérifie si le nombre est pair en utilisant l'opérateur modulo (\%) pour voir s'il n'y a pas de reste lorsque l'on divise par 2.
    print(evens\_squared) : Cette ligne affiche le dictionnaire evens\_squared sur la console.
        \end{solution}
        

        \question
        Créer un dictionnaire de mots et de leur longueur à partir d'une phrase

Exemple de sortie

Enter String : Python is awesome (Python est génial)

\{'Python' : 6, 'is' : 2, 'awesome' : 7\}
        \par
        \begin{solution}
            \renewcommand{\nomfichier}{q662.py}
            \pythonfile{\chemincode \nomfichier}[][\nomfichier]
            Ce code Python crée un dictionnaire nommé word\_lengths dont les clés sont les mots d'une phrase et les valeurs sont les longueurs de ces mots. Voici comment fonctionne ce code :

    sentence = "Python is awesome" : Cette ligne initialise une variable de chaîne nommée phrase avec la phrase que vous souhaitez analyser.
    word\_lengths = \{word : len(word) for word in sentence.split()\} : Cette ligne utilise une compréhension de dictionnaire pour créer le dictionnaire word\_lengths. Elle divise la phrase en mots à l'aide de .split() et parcourt chaque mot de la phrase. Pour chaque mot, elle assigne une paire clé-valeur au dictionnaire. La clé est le mot lui-même (word), et la valeur est la longueur de ce mot (len(word)).
        pour mot dans phrase.split() : Cette partie du code parcourt chaque mot de la phrase après l'avoir divisée en mots.
    print(sentence) : Cette ligne imprime la phrase originale sur la console.
    print(word\_lengths) : Cette ligne affiche le dictionnaire word\_lengths, qui contient les longueurs des mots de la phrase originale.
        \end{solution}
        

        \question
        Créer un dictionnaire de caractères minuscules à partir d'une chaîne de caractères

Exemple de sortie

Bonjour à tous

\{'H' : 'h', 'e' : 'e', 'l' : 'l', 'o' : 'o', 'W' : 'w', 'r' : 'r', 'd' : 'd'\}
        \par
        \begin{solution}
            \renewcommand{\nomfichier}{q663.py}
            \pythonfile{\chemincode \nomfichier}[][\nomfichier]
            Ce code Python crée un dictionnaire nommé lowercase\_chars dont les clés sont les caractères alphabétiques d'un texte donné et les valeurs sont les versions minuscules correspondantes de ces caractères. Voici comment fonctionne ce code :

    text = "Hello World" : Cette ligne initialise une variable de chaîne nommée texte avec le texte que vous souhaitez analyser.
    lowercase\_chars = \{char : char.lower() for char in text if char.isalpha()\} : Cette ligne utilise une compréhension de dictionnaire pour créer le dictionnaire lowercase\_chars. Elle parcourt chaque caractère char du texte et vérifie s'il s'agit d'un caractère alphabétique (en utilisant char.isalpha()). Pour chaque caractère alphabétique, il attribue une paire clé-valeur au dictionnaire. La clé est le caractère lui-même (char) et la valeur est la version minuscule de ce caractère (char.lower()).
        pour char dans text : Cette partie du code parcourt chaque caractère du texte.
        if char.isalpha() : Cette partie du code vérifie si le caractère est alphabétique.
    print(text) : Cette ligne imprime le texte original sur la console.
    print(lowercase\_chars) : Cette ligne affiche le dictionnaire lowercase\_chars, qui contient les versions minuscules des caractères alphabétiques du texte original.
        \end{solution}
        

        \question
        Créer un dictionnaire des nombres et de leurs cubes

Exemple de résultat

[1, 2, 3, 4, 5]

\{1 : 1, 2 : 8, 3 : 27, 4 : 64, 5 : 125\}
        \par
        \begin{solution}
            \renewcommand{\nomfichier}{q664.py}
            \pythonfile{\chemincode \nomfichier}[][\nomfichier]
            Ce code Python crée un dictionnaire nommé cubes où les clés sont des nombres d'une liste, et les valeurs sont les cubes de ces nombres. Voici comment fonctionne le code :

    nombres = [1, 2, 3, 4, 5] : Cette ligne initialise une liste nommée numbers avec une séquence de nombres que vous souhaitez analyser.
    cubes = \{x : x ** 3 for x in numbers\} : Cette ligne utilise une compréhension de dictionnaire pour créer le dictionnaire cubes. Elle parcourt chaque nombre x dans la liste des nombres et affecte une paire clé-valeur au dictionnaire. La clé est le nombre lui-même (x) et la valeur est le cube de ce nombre (x ** 3).
        pour x dans nombres : Cette partie du code parcourt chaque nombre de la liste des nombres.
    print(numbers) : Cette ligne imprime la liste originale des nombres sur la console.
    print(cubes) : Cette ligne imprime le dictionnaire cubes, qui contient les cubes des nombres de la liste originale.
        \end{solution}
        

        \question
        Créer un dictionnaire des nombres et de leurs carrés, à l'exclusion des nombres impairs

Exemple de résultat

[1, 2, 3, 4, 5, 6, 7, 8, 9, 10]

\{2 : 4, 4 : 16, 6 : 36, 8 : 64, 10 : 100\}
        \par
        \begin{solution}
            \renewcommand{\nomfichier}{q665.py}
            \pythonfile{\chemincode \nomfichier}[][\nomfichier]
            Ce code Python crée un dictionnaire nommé even\_squares dont les clés sont les nombres pairs d'une liste et les valeurs sont les carrés de ces nombres pairs. Voici comment fonctionne ce code :

    numbers = [1, 2, 3, 4, 5, 6, 7, 8, 9, 10] : Cette ligne initialise une liste nommée numbers avec une séquence de nombres que vous souhaitez analyser.
    even\_squares = \{x : x ** 2 for x in numbers if x \% 2 == 0\} : Cette ligne utilise une compréhension de dictionnaire pour créer le dictionnaire even\_squares. Elle parcourt chaque nombre x dans la liste des nombres et vérifie si le nombre est pair (en utilisant x \% 2 == 0). Pour chaque nombre pair, il attribue une paire clé-valeur au dictionnaire. La clé est le nombre pair lui-même (x) et la valeur est le carré de ce nombre (x ** 2).
        pour x dans nombres : Cette partie du code parcourt chaque nombre de la liste des nombres.
        if x \% 2 == 0 : Cette partie du code vérifie si le nombre est pair (c'est-à-dire si son reste après division par 2 est égal à 0).
    print(even\_squares) : Cette ligne affiche le dictionnaire even\_squares, qui contient les carrés des nombres pairs de la liste originale.
        \end{solution}
        

        \question
        Créer un dictionnaire des nombres et de leur statut premier

Exemple de résultat

[1, 2, 3, 4, 5, 6, 7, 8, 9, 10]

\{1 : Faux, 2 : Vrai, 3 : Vrai, 4 : Faux, 5 : Vrai, 6 : Faux, 7 : Vrai, 8 : Faux, 9 : Faux, 10 : Faux\}
        \par
        \begin{solution}
            \renewcommand{\nomfichier}{q666.py}
            \pythonfile{\chemincode \nomfichier}[][\nomfichier]
            Ce code Python définit une fonction is\_prime(n) qui vérifie si un nombre donné n est premier. Il utilise ensuite une compréhension de dictionnaire pour créer un dictionnaire nommé prime\_status où les clés sont des nombres d'une liste et les valeurs sont des valeurs booléennes indiquant si chaque nombre est premier ou non. Voici comment fonctionne le code :

    def is\_prime(n) : Ceci définit une fonction is\_prime qui prend un seul argument n et vérifie s'il s'agit d'un nombre premier. La fonction renvoie True si n est premier et False dans le cas contraire. Elle utilise un algorithme standard de vérification de la primalité, testant les diviseurs jusqu'à la racine carrée de n.
    numbers = [1, 2, 3, 4, 5, 6, 7, 8, 9, 10] : Cette ligne initialise une liste nommée numbers avec une séquence de nombres que vous souhaitez analyser.
    prime\_status = \{x : is\_prime(x) for x in numbers\} : Cette ligne utilise une compréhension de dictionnaire pour créer le dictionnaire prime\_status. Elle parcourt chaque nombre x dans la liste des nombres et affecte une paire clé-valeur au dictionnaire. La clé est le nombre lui-même (x) et la valeur est le résultat de l'appel à la fonction is\_prime sur ce nombre, qui détermine s'il est premier ou non.
        pour x dans nombres : Cette partie du code parcourt chaque nombre de la liste des nombres.
        {x : is\_prime(x) for x in numbers} : Elle utilise la compréhension d'un dictionnaire pour créer des paires clé-valeur où la clé est le nombre x et la valeur est le résultat de la fonction is\_prime pour ce nombre.
    print(nombres) : Cette ligne imprime la liste originale des nombres sur la console.
    print(prime\_status) : Cette ligne affiche le dictionnaire prime\_status, qui contient le statut premier (Vrai ou Faux) de chaque nombre de la liste originale.
        \end{solution}
        

        \question
        Créer un dictionnaire de caractères et de leurs valeurs ASCII à partir d'une chaîne de caractères

Exemple de sortie

Tuteur Joes

\{'T' : 84, 'u' : 117, 't' : 116, 'o' : 111, 'r' : 114, ' ' : 32, 'J' : 74, 'e' : 101, 's' : 115\}
        \par
        \begin{solution}
            \renewcommand{\nomfichier}{q667.py}
            \pythonfile{\chemincode \nomfichier}[][\nomfichier]
            Ce code Python prend une chaîne de texte et crée un dictionnaire ascii\_values où les clés sont des caractères de la chaîne, et les valeurs sont les valeurs ASCII correspondantes. Voici comment fonctionne ce code :

    text = "Tutor Joes" : Cette ligne initialise une variable texte avec la valeur "Tutor Joes".
    ascii\_values = \{char : ord(char) for char in text\} : Cette ligne utilise une compréhension de dictionnaire pour créer le dictionnaire ascii\_values. Elle parcourt chaque caractère char de la chaîne de texte et affecte une paire clé-valeur au dictionnaire. La clé est le caractère lui-même (char) et la valeur est la valeur ASCII de ce caractère, obtenue à l'aide de la fonction ord(char).
        pour char dans text : Cette partie du code parcourt chaque caractère de la chaîne de texte.
        {char : ord(char) for char in text} : Elle utilise la compréhension d'un dictionnaire pour créer des paires clé-valeur où la clé est le caractère char et la valeur est sa valeur ASCII correspondante.
    print(text) : Cette ligne imprime la chaîne de texte originale sur la console.
    print(ascii\_values) : Cette ligne affiche le dictionnaire ascii\_values, qui contient les caractères de la chaîne originale comme clés et leurs valeurs ASCII comme valeurs.
        \end{solution}
        

        \question
        Créer un dictionnaire de mots et de leurs voyelles à partir d'une liste de chaînes de caractères.

Exemple de résultat

['apple', 'banana', 'cherry']

\{'apple' : 2, 'banana' : 3, 'cherry' : 1\}
        \par
        \begin{solution}
            \renewcommand{\nomfichier}{q668.py}
            \pythonfile{\chemincode \nomfichier}[][\nomfichier]
            Ce code Python prend une liste de mots et crée un dictionnaire vowel\_counts dont les clés sont les mots et les valeurs sont les nombres de voyelles (a, e, i, o, u) dans chaque mot. Voici comment fonctionne le code :

    words = ['apple', 'banana', 'cherry'] : Cette ligne initialise une liste nommée words avec trois mots.
    vowel\_counts = \{word : sum(1 for char in word if char.lower() in 'aeiou') for word in words\} : Cette ligne utilise une compréhension de dictionnaire pour créer le dictionnaire vowel\_counts. Elle parcourt chaque mot de la liste des mots et assigne une paire clé-valeur au dictionnaire. La clé est le mot lui-même (word), et la valeur est le résultat de l'expression suivante :
        sum(1 for char in word if char.lower() in 'aeiou') : Cette partie du code compte le nombre de voyelles dans chaque mot en parcourant chaque caractère char dans le mot. Si char est une voyelle (quelle que soit la casse), il ajoute 1 au compte. La fonction sum calcule le nombre total de voyelles dans le mot.
        pour mot dans mots : Cette partie du code parcourt chaque mot de la liste des mots.
    print(words) : Cette ligne imprime la liste de mots originale sur la console.
    print(vowel\_counts) : Cette ligne imprime le dictionnaire vowel\_counts, qui contient les mots comme clés et le nombre de voyelles dans chaque mot comme valeurs.
        \end{solution}
        

        \question
        Créer un dictionnaire de mots avec les lettres triées

Exemple de résultat

['python', 'programmation', 'langage']

\{'python' : 'hnopty', 'programming' : 'aggimmnoprr', 'language' : 'aaegglnu'\}
        \par
        \begin{solution}
            \renewcommand{\nomfichier}{q669.py}
            \pythonfile{\chemincode \nomfichier}[][\nomfichier]
            Ce code Python prend une liste de mots et crée un dictionnaire sorted\_letters où les clés sont des mots, et les valeurs sont les mots avec leurs lettres triées dans l'ordre alphabétique. Voici comment fonctionne ce code :

    mots = ['python', 'programmation', 'langage'] : Cette ligne initialise une liste nommée words avec trois mots.
    sorted\_letters = \{word : ''.join(sorted(word)) for word in words\} : Cette ligne utilise une compréhension de dictionnaire pour créer le dictionnaire sorted\_letters. Elle parcourt chaque mot de la liste des mots et affecte une paire clé-valeur au dictionnaire. La clé est le mot lui-même (word), et la valeur est le résultat de l'expression suivante :
        ''.join(sorted(word)) : Cette partie du code trie les lettres de chaque mot par ordre alphabétique et les réunit en une chaîne. La fonction sorted est utilisée pour trier les caractères du mot, et join combine les caractères triés en une seule chaîne.
        pour mot dans mots : Cette partie du code parcourt chaque mot de la liste des mots.
    print(words) : Cette ligne imprime la liste de mots originale sur la console.
    print(lettres\_triées) : Cette ligne affiche le dictionnaire sorted\_letters, qui contient les mots comme clés et les mots avec leurs lettres triées par ordre alphabétique comme valeurs.
        \end{solution}
        

        \question
        Créer un dictionnaire des mots et de leur longueur, mais seulement pour les mots de plus de 5 lettres.

Exemple de résultat

['apple', 'banana', 'cherry', 'date']

\{'banane' : 6, 'cerise' : 6\}
        \par
        \begin{solution}
            \renewcommand{\nomfichier}{q670.py}
            \pythonfile{\chemincode \nomfichier}[][\nomfichier]
            Ce code Python prend une liste de mots et crée un dictionnaire long\_word\_lengths où les clés sont les mots dont la longueur est supérieure à 5 caractères, et les valeurs sont les longueurs de ces mots. Voici comment fonctionne ce code :

    words = ['apple', 'banana', 'cherry', 'date'] : Cette ligne initialise une liste nommée words avec quatre mots.
    long\_word\_lengths = \{word : len(word) for word in words if len(word) > 5\} : Cette ligne utilise une compréhension de dictionnaire pour créer le dictionnaire long\_word\_lengths. Elle parcourt chaque mot de la liste des mots et affecte une paire clé-valeur au dictionnaire uniquement si la longueur du mot est supérieure à 5 caractères. La clé est le mot lui-même (word) et la valeur est la longueur du mot (calculée à l'aide de len(word)).
        pour mot dans mots : Cette partie du code parcourt chaque mot de la liste des mots.
        if len(word) > 5 : Cette partie du code vérifie si la longueur du mot courant est supérieure à 5.
    print(words) : Cette ligne imprime la liste de mots originale sur la console.
    print(long\_word\_lengths) : Cette ligne affiche le dictionnaire long\_word\_lengths, qui contient les mots dont la longueur est supérieure à 5 comme clés et leurs longueurs respectives comme valeurs.
        \end{solution}
        

        \question
        Créer un dictionnaire des caractères et de leur fréquence dans une chaîne de caractères

Exemple de sortie

bonjour le monde

\{'l' : 3, ' ' : 1, 'r' : 1, 'h' : 1, 'd' : 1, 'o' : 2, 'e' : 1, 'w' : 1\}
        \par
        \begin{solution}
            \renewcommand{\nomfichier}{q671.py}
            \pythonfile{\chemincode \nomfichier}[][\nomfichier]
            Ce code Python prend une chaîne de texte et crée un dictionnaire char\_frequency où les clés sont des caractères uniques dans la chaîne, et les valeurs sont les fréquences de ces caractères dans la chaîne. Voici comment fonctionne le code :

    text = "hello world" : Cette ligne initialise la chaîne de caractères text avec le texte "hello world".
    char\_frequency = \{char : text.count(char) for char in set(text)\} : Cette ligne utilise une compréhension de dictionnaire pour créer le dictionnaire char\_frequency. Elle effectue les opérations suivantes :
        set(text) : Cette partie du code crée un ensemble de caractères uniques dans la chaîne de texte. L'utilisation d'un ensemble permet de s'assurer que chaque caractère n'est inclus qu'une seule fois, ce qui élimine les doublons.
        \{char : text.count(char) for char in set(text)\} : Cette partie du code parcourt chaque caractère unique char dans le jeu de caractères uniques et attribue une paire clé-valeur au dictionnaire. La clé est le caractère lui-même (char) et la valeur est le résultat de text.count(char), qui compte le nombre d'occurrences de ce caractère dans la chaîne de texte.
    print(text) : Cette ligne imprime la chaîne de texte originale sur la console.
    print(char\_frequency) : Cette ligne affiche le dictionnaire char\_frequency, qui contient des caractères uniques comme clés et leurs fréquences comme valeurs.
        \end{solution}
        

        \question
        Créer un dictionnaire de mots et de leurs formes inversées

Exemple de résultat

['apple', 'banana', 'cherry']

\{'apple' : 'elppa', 'banana' : 'ananab', 'cherry' : 'yrrehc'\}.
        \par
        \begin{solution}
            \renewcommand{\nomfichier}{q672.py}
            \pythonfile{\chemincode \nomfichier}[][\nomfichier]
            Ce code Python prend une liste de mots et crée un dictionnaire reversed\_words dont les clés sont les mots originaux et les valeurs les versions inversées de ces mots. Voici comment fonctionne ce code :

    mots = ['pomme', 'banane', 'cerise'] : Cette ligne initialise une liste nommée words avec trois mots.
    reversed\_words = \{word : word[::-1] for word in words\} : Cette ligne utilise une compréhension de dictionnaire pour créer le dictionnaire mots\_inversés. Elle parcourt chaque mot de la liste des mots et affecte une paire clé-valeur au dictionnaire. La clé est le mot original (word), et la valeur est le mot inversé en utilisant le découpage (word[::-1]).
        pour mot dans mots : Cette partie du code parcourt chaque mot de la liste des mots.
        mot[::-1] : Cette partie du code découpe le mot avec un pas de -1, inversant ainsi l'ordre des caractères dans le mot.
    print(mots) : Cette ligne imprime la liste de mots originale sur la console.
    print(mots\_inversés) : Cette ligne affiche le dictionnaire mots\_inversés, qui contient les mots originaux comme clés et leurs versions inversées comme valeurs.
        \end{solution}
        

        \question
        Créer un dictionnaire des nombres et de leurs carrés, mais uniquement pour les nombres pairs

Exemple de résultat

[1, 2, 3, 4, 5, 6, 7, 8, 9, 10]

\{2 : 4, 4 : 16, 6 : 36, 8 : 64, 10 : 100\}
        \par
        \begin{solution}
            \renewcommand{\nomfichier}{q673.py}
            \pythonfile{\chemincode \nomfichier}[][\nomfichier]
            Ce code Python prend une liste de nombres et crée un dictionnaire even\_squares dont les clés sont les nombres pairs de la liste originale, et les valeurs sont les carrés de ces nombres pairs. Voici comment fonctionne ce code :

    nombres = [1, 2, 3, 4, 5, 6, 7, 8, 9, 10] : Cette ligne initialise une liste nommée numbers avec dix valeurs entières.
    even\_squares = \{x : x ** 2 for x in numbers if x \% 2 == 0\} : Cette ligne utilise une compréhension de dictionnaire pour créer le dictionnaire even\_squares. Elle parcourt chaque nombre x dans la liste des nombres et affecte une paire clé-valeur au dictionnaire si et seulement si le nombre est pair (c'est-à-dire si x \% 2 == 0 est Vrai). La clé est le nombre pair lui-même (x) et la valeur est le carré de ce nombre pair (x ** 2).
        pour x dans nombres : Cette partie du code parcourt chaque nombre de la liste des nombres.
        if x \% 2 == 0 : Cette partie du code vérifie si le nombre est pair en testant s'il est divisible par 2 sans reste.
    print(numbers) : Cette ligne imprime la liste originale des nombres sur la console.
    print(even\_squares) : Cette ligne affiche le dictionnaire even\_squares, qui contient les nombres pairs comme clés et leurs valeurs au carré comme valeurs.
        \end{solution}
        

        \question
        Créer un dictionnaire des nombres et de leurs facteurs

Exemple de résultat

[1, 2, 3, 4, 5]

\{1 : [1], 2 : [1, 2], 3 : [1, 3], 4 : [1, 2, 4], 5 : [1, 5]\}
        \par
        \begin{solution}
            \renewcommand{\nomfichier}{q674.py}
            \pythonfile{\chemincode \nomfichier}[][\nomfichier]
            Ce code Python prend une liste de nombres et crée un dictionnaire factors où les clés sont les nombres de la liste originale, et les valeurs sont des listes de leurs facteurs. Voici comment fonctionne le code :

    nombres = [1, 2, 3, 4, 5] : Cette ligne initialise une liste nommée numbers avec cinq valeurs entières.
    facteurs = \{x : [i for i in range(1, x + 1) if x \% i == 0] for x in numbers\} : Cette ligne utilise une compréhension de dictionnaire pour créer le dictionnaire des facteurs. Elle parcourt chaque nombre x dans la liste des nombres et affecte une paire clé-valeur au dictionnaire. La clé est le nombre lui-même (x) et la valeur est une liste de ses facteurs.
        pour x dans nombres : Cette partie du code parcourt chaque nombre de la liste des nombres.
        x : [i for i in range(1, x + 1) if x \% i == 0] : Cette partie du code utilise une compréhension de liste pour générer une liste de facteurs pour chaque nombre. Elle parcourt les nombres de 1 à x (inclus) et n'inclut que les nombres pour lesquels x \% i == 0 est Vrai. Cette condition vérifie si i est un facteur de x.
    print(nombres) : Cette ligne imprime la liste originale des nombres sur la console.
    print(facteurs) : Cette ligne affiche le dictionnaire des facteurs, qui contient des nombres comme clés et des listes de leurs facteurs comme valeurs.
        \end{solution}
        

        \question
        Créer un dictionnaire de mots et de leurs majuscules

Exemple de sortie

['apple', 'banana', 'cherry']

\{'apple' : 'APPLE', 'banana' : 'BANANA', 'cherry' : 'CHERRY'\}.
        \par
        \begin{solution}
            \renewcommand{\nomfichier}{q675.py}
            \pythonfile{\chemincode \nomfichier}[][\nomfichier]
            Ce code Python prend une liste de mots et crée un dictionnaire uppercase\_words dont les clés sont les mots de la liste originale et les valeurs sont les versions en majuscules de ces mots. Voici comment fonctionne ce code :

    mots = ['pomme', 'banane', 'cerise'] : Cette ligne initialise une liste nommée words avec trois valeurs de chaîne.
    uppercase\_words = \{word : word.upper() for word in words\} : Cette ligne utilise une compréhension de dictionnaire pour créer le dictionnaire uppercase\_words. Elle parcourt chaque mot de la liste des mots et affecte une paire clé-valeur au dictionnaire. La clé est le mot original (word) et la valeur est la version en majuscules du mot (word.upper()).
        pour mot dans mots : Cette partie du code parcourt chaque mot de la liste des mots.
        word : word.upper() : Cette partie du code crée la paire clé-valeur. La clé est le mot original et la valeur est la version en majuscules du mot.
    print(words) : Cette ligne imprime la liste des mots originaux sur la console.
    print(mots\_en\_majuscules) : Cette ligne affiche le dictionnaire uppercase\_words, qui contient des mots comme clés et leurs versions en majuscules comme valeurs.
        \end{solution}
        

        \question
        Créer un dictionnaire des nombres et de leur parité (paire ou impaire)

Exemple de résultat

[1, 2, 3, 4, 5, 6, 7, 8, 9, 10]

\{1 : 'odd', 2 : 'even', 3 : 'odd', 4 : 'even', 5 : 'odd', 6 : 'even', 7 : 'odd', 8 : 'even', 9 : 'odd', 10 : 'even'\}
        \par
        \begin{solution}
            \renewcommand{\nomfichier}{q676.py}
            \pythonfile{\chemincode \nomfichier}[][\nomfichier]
            Ce code Python prend une liste de nombres et crée un dictionnaire de parité où les clés sont les nombres de la liste originale, et les valeurs indiquent si chaque nombre est "pair" ou "impair". Voici comment fonctionne ce code :

    nombres = [1, 2, 3, 4, 5, 6, 7, 8, 9, 10] : Cette ligne initialise une liste nommée nombres avec dix valeurs entières.
    parité = \{x : 'even' if x \% 2 == 0 else 'odd' for x in numbers\} : Cette ligne utilise une compréhension de dictionnaire pour créer le dictionnaire de parité. Elle parcourt chaque nombre x dans la liste des nombres et affecte une paire clé-valeur au dictionnaire. La clé est le nombre lui-même (x) et la valeur est déterminée à l'aide d'une expression conditionnelle :
        Si x \% 2 == 0, il est considéré comme "pair". Sinon, il est considéré comme "impair".
        pour x dans nombres : Cette partie du code parcourt chaque nombre de la liste des nombres.
        x : "pair" si x \% 2 == 0 sinon "impair" : Cette partie du code crée la paire clé-valeur. La clé est le nombre x, et la valeur est soit "pair", soit "impair", en fonction du résultat de l'expression conditionnelle.
    print(nombres) : Cette ligne imprime la liste originale des nombres sur la console.
    print(parity) : Cette ligne affiche le dictionnaire de parité, qui contient des nombres comme clés et leur parité (paire ou impaire) comme valeurs.
        \end{solution}
        

        \question
        Créer un dictionnaire des nombres et de leurs représentations binaires

Exemple de sortie

[1, 2, 3, 4, 5]

\{1 : '0b1', 2 : '0b10', 3 : '0b11', 4 : '0b100', 5 : '0b101'\}
        \par
        \begin{solution}
            \renewcommand{\nomfichier}{q677.py}
            \pythonfile{\chemincode \nomfichier}[][\nomfichier]
            Ce code Python prend une liste de nombres et crée un dictionnaire binary\_representations dont les clés sont les nombres de la liste originale et les valeurs sont leurs représentations binaires sous forme de chaînes de caractères. Voici comment fonctionne ce code :

    nombres = [1, 2, 3, 4, 5] : Cette ligne initialise une liste nommée numbers avec cinq valeurs entières.
    binary\_representations = \{x : bin(x) for x in numbers\} : Cette ligne utilise une compréhension de dictionnaire pour créer le dictionnaire binary\_representations. Elle parcourt chaque nombre x de la liste des nombres et affecte une paire clé-valeur au dictionnaire. La clé est le nombre lui-même (x) et la valeur est sa représentation binaire sous forme de chaîne, obtenue à l'aide de la fonction bin.
        pour x dans nombres : Cette partie du code parcourt chaque nombre de la liste des nombres.
        x : bin(x) : Cette partie du code crée la paire clé-valeur. La clé est le nombre x et la valeur est sa représentation binaire sous forme de chaîne de caractères.
    print(nombres) : Cette ligne imprime la liste originale des nombres sur la console.
    print(binary\_representations) : Cette ligne affiche le dictionnaire binary\_representations, qui contient des nombres comme clés et leurs représentations binaires comme valeurs.
        \end{solution}
        

        \question
        Créer un dictionnaire des nombres et de leurs factoriels

Exemple de résultat

[1, 2, 3, 4, 5]

\{1 : 1, 2 : 2, 3 : 6, 4 : 24, 5 : 120\}
        \par
        \begin{solution}
            \renewcommand{\nomfichier}{q678.py}
            \pythonfile{\chemincode \nomfichier}[][\nomfichier]
            Ce code Python crée un dictionnaire factorials dont les clés sont des nombres issus d'une liste et les valeurs sont leurs valeurs factorielles calculées à l'aide de la fonction math.factorial. Voici comment fonctionne le code :

    nombres = [1, 2, 3, 4, 5] : Cette ligne initialise une liste nommée numbers avec cinq valeurs entières.
    factorials = \{x : math.factorial(x) for x in numbers\} : Cette ligne utilise une compréhension de dictionnaire pour créer le dictionnaire factorials. Elle parcourt chaque nombre x dans la liste des nombres et affecte une paire clé-valeur au dictionnaire. La clé est le nombre lui-même (x) et la valeur est sa factorielle calculée à l'aide de la fonction math.factorial.
        pour x dans nombres : Cette partie du code parcourt chaque nombre de la liste des nombres.
        x : math.factorial(x) : Cette partie du code crée la paire clé-valeur. La clé est le nombre x et la valeur est sa valeur factorielle calculée à l'aide de math.factorial(x).
    print(nombres) : Cette ligne imprime la liste originale des nombres sur la console.
    print(factorials) : Cette ligne affiche le dictionnaire factorials, qui contient des nombres comme clés et leurs valeurs factorielles comme valeurs.
        \end{solution}
        

        \question
        Créer un dictionnaire des nombres et de leurs carrés, mais seulement pour les multiples de 3

Exemple de résultat

[1, 2, 3, 4, 5, 6, 7, 8, 9, 10]

\{3 : 9, 6 : 36, 9 : 81\}
        \par
        \begin{solution}
            \renewcommand{\nomfichier}{q679.py}
            \pythonfile{\chemincode \nomfichier}[][\nomfichier]
            Ce code Python crée un dictionnaire multiples\_of\_3\_squares où les clés sont des nombres d'une liste, mais seulement pour les nombres qui sont des multiples de 3, et les valeurs sont les carrés de ces nombres. Voici comment fonctionne le code :

    nombres = [1, 2, 3, 4, 5, 6, 7, 8, 9, 10] : Cette ligne initialise une liste nommée nombres avec dix valeurs entières.
    multiples\_of\_3\_squares = \{x : x ** 2 for x in numbers if x \% 3 == 0\} : Cette ligne utilise une compréhension de dictionnaire pour créer le dictionnaire multiples\_of\_3\_squares. Elle parcourt chaque nombre x dans la liste des nombres et affecte une paire clé-valeur au dictionnaire uniquement si le nombre est un multiple de 3 (c'est-à-dire lorsque x \% 3 == 0). La clé est le nombre lui-même (x) et la valeur est le carré de ce nombre (x ** 2).
        pour x dans nombres : Cette partie du code parcourt chaque nombre de la liste des nombres.
        if x \% 3 == 0 : cette partie du code vérifie si le nombre x est un multiple de 3.
        x : x ** 2 : cette partie du code crée la paire clé-valeur. La clé est le nombre x et la valeur est son carré, calculé comme x ** 2.
    print(nombres) : Cette ligne imprime la liste originale des nombres sur la console.
    print(multiples\_de\_3\_squares) : Cette ligne affiche le dictionnaire multiples\_de\_3\_squares, qui contient des nombres (multiples de 3) comme clés et leurs valeurs au carré comme valeurs.
        \end{solution}
        

        \question
        Créer un dictionnaire des nombres et de leurs puissances de 2

Exemple de sortie

[1, 2, 3, 4, 5]

\{1 : 2, 2 : 4, 3 : 8, 4 : 16, 5 : 32\}
        \par
        \begin{solution}
            \renewcommand{\nomfichier}{q680.py}
            \pythonfile{\chemincode \nomfichier}[][\nomfichier]
            Ce code Python crée un dictionnaire powers\_of\_2, où les clés sont des nombres d'une liste, et les valeurs sont 2 élevés à la puissance de chaque nombre. Voici comment fonctionne le code :

    nombres = [1, 2, 3, 4, 5] : Cette ligne initialise une liste nommée numbers avec cinq valeurs entières.
    puissances\_de\_2 = \{x : 2 ** x for x in numbers\} : Cette ligne utilise une compréhension de dictionnaire pour créer le dictionnaire powers\_of\_2. Elle parcourt chaque nombre x dans la liste des nombres et affecte une paire clé-valeur au dictionnaire. La clé est le nombre lui-même (x), et la valeur est calculée en élevant 2 à la puissance de ce nombre (2 ** x).
        pour x dans nombres : Cette partie du code parcourt chaque nombre de la liste des nombres.
        x : 2 ** x : Cette partie du code crée la paire clé-valeur. La clé est le nombre x, et la valeur est calculée comme 2 ** x, ce qui élève 2 à la puissance du nombre x.
    print(numbers) : Cette ligne imprime la liste originale des nombres sur la console.
    print(puissances\_de\_2) : Cette ligne affiche le dictionnaire powers\_of\_2, qui contient des nombres comme clés et leurs puissances de 2 correspondantes comme valeurs.
        \end{solution}
        

        \question
        Créer un dictionnaire de mots et de leurs voyelles, en excluant les mots sans voyelles.

Exemple de résultat

['apple', 'banana', 'cherry', 'dog', 'ct']

\{'apple' : 2, 'banana' : 3, 'cherry' : 1, 'dog' : 1, 'cat' : 1\}
        \par
        \begin{solution}
            \renewcommand{\nomfichier}{q681.py}
            \pythonfile{\chemincode \nomfichier}[][\nomfichier]
            Ce code Python définit un dictionnaire vowel\_counts, dont les clés sont des mots d'une liste contenant au moins une voyelle (a, e, i, o, u), et dont les valeurs sont les nombres de voyelles dans chaque mot. Voici comment fonctionne le code :

    words = ['apple', 'banana', 'cherry', 'dog', 'ct'] : Cette ligne initialise une liste nommée words avec cinq valeurs de chaîne, dont certaines contiennent des voyelles.
    vowel\_counts = \{word : sum(1 for char in word if char.lower() in 'aeiou') for word in words if any(char.lower() in 'aeiou' for char in word)\} : Cette ligne utilise une compréhension du dictionnaire pour créer le dictionnaire vowel\_counts. Elle parcourt chaque mot de la liste des mots, et pour chaque mot, elle vérifie s'il y a au moins une voyelle (n'importe quel caractère voyelle) dans ce mot.
        pour mot dans mots : Cette partie du code parcourt chaque mot de la liste des mots.
        any(char.lower() in 'aeiou' for char in word) : Cette partie vérifie si l'un des caractères du mot est une voyelle minuscule ("a", "e", "i", "o", "u"). La partie char.lower() garantit que les voyelles minuscules et majuscules sont comptées.
        word : sum(1 for char in word if char.lower() in 'aeiou') : Si le mot contient au moins une voyelle, une paire clé-valeur est créée dans le dictionnaire. La clé est le mot lui-même (word), et la valeur est calculée en additionnant 1 pour chaque caractère du mot qui est une voyelle. La fonction char.lower() de la partie 'aeiou' vérifie si le caractère est une voyelle minuscule.
    print(words) : Cette ligne affiche la liste originale des mots sur la console.
    print(vowel\_counts) : Cette ligne affiche le dictionnaire vowel\_counts, qui contient les mots comportant au moins une voyelle et leur nombre de voyelles.
        \end{solution}
        

        \question
        Créer un dictionnaire de chaînes de caractères en remplaçant les voyelles par des traits de soulignement

Exemple de résultat

['apple', 'banana', 'cherry']

\{'apple' : '\_ppl\_', 'banana' : 'b\_n\_n\_', 'cherry' : 'ch\_rry'\}
        \par
        \begin{solution}
            \renewcommand{\nomfichier}{q682.py}
            \pythonfile{\chemincode \nomfichier}[][\nomfichier]
            Le code Python définit un dictionnaire underscored\_words, dont les clés sont des mots d'une liste et les valeurs sont les mêmes mots dont les voyelles ont été remplacées par des traits de soulignement. Voici comment fonctionne le code :

    words = ['apple', 'banana', 'cherry'] : Cette ligne initialise une liste nommée words avec trois valeurs de chaîne.
    underscored\_words = \{word : ''.join(['\_' if char.lower() in 'aeiou' else char for char in word]) for word in words\} : Cette ligne utilise une compréhension de dictionnaire pour créer le dictionnaire underscored\_words. Elle parcourt chaque mot de la liste des mots et crée une paire clé-valeur pour chaque mot.
        for word in words : Cette partie du code parcourt chaque mot de la liste des mots.
        ''.join(['\_' if char.lower() in 'aeiou' else char for char in word]) : Pour chaque mot, il crée une version modifiée du mot dont les voyelles sont remplacées par des traits de soulignement ("\_"). Pour ce faire, il parcourt chaque caractère du mot (for char in word) et utilise une expression conditionnelle pour vérifier si le caractère est une voyelle minuscule ('a', 'e', 'i', 'o', 'u'). S'il s'agit d'une voyelle, il la remplace par un trait de soulignement ; sinon, il conserve le caractère tel quel. La méthode join est utilisée pour concaténer les caractères modifiés en une seule chaîne.
    print(words) : Cette ligne affiche la liste originale des mots sur la console.
    print(mots\_soulignés) : Cette ligne affiche le dictionnaire underscored\_words, qui contient les mots originaux comme clés et leurs versions remplacées par des voyelles comme valeurs.
        \end{solution}
        

        \question
        Créer un dictionnaire de mots et de leur longueur, mais uniquement pour les mots commençant par "a".

Exemple de résultat

['apple', 'banana', 'cherry', 'date']

\{'apple' : 5\}
        \par
        \begin{solution}
            \renewcommand{\nomfichier}{q683.py}
            \pythonfile{\chemincode \nomfichier}[][\nomfichier]
            Le code Python que vous avez fourni crée un dictionnaire a\_word\_lengths, dont les clés sont les mots d'une liste commençant par la lettre "a", et les valeurs sont les longueurs de ces mots. Voici comment fonctionne le code :

    words = ['apple', 'banana', 'cherry', 'date'] : Cette ligne initialise une liste nommée words avec quatre valeurs de chaîne.
    a\_word\_lengths = \{word : len(word) for word in words if word.startswith('a')\} : Cette ligne utilise une compréhension de dictionnaire pour créer le dictionnaire a\_word\_lengths. Elle parcourt chaque mot de la liste des mots et crée une paire clé-valeur pour chaque mot commençant par la lettre "a".
        pour mot dans mots : Cette partie du code parcourt chaque mot de la liste des mots.
        'a\_word\_lengths' = \{word : len(word) for word in words if word.startswith('a')\} : Pour chaque mot, il est vérifié si le mot commence par la lettre "a" en utilisant la condition word.startswith('a'). Si la condition est vraie, le mot est inclus dans le dictionnaire. La clé est le mot lui-même, et la valeur est la longueur du mot, calculée à l'aide de la fonction len(word).
    print(words) : Cette ligne affiche la liste originale des mots sur la console.
    print(a\_word\_lengths) : Cette ligne affiche le dictionnaire a\_word\_lengths, qui contient les mots commençant par "a" comme clés et leurs longueurs correspondantes comme valeurs.
        \end{solution}
        

        \question
        Créer un dictionnaire de chaînes de caractères en répétant chaque mot trois fois

Exemple de résultat

['apple', 'banana', 'cherry']

\{'apple' : 'appleappleapple', 'banana' : 'bananabanabanana', 'cherry' : 'cherrycherrycherry'\}.
        \par
        \begin{solution}
            \renewcommand{\nomfichier}{q684.py}
            \pythonfile{\chemincode \nomfichier}[][\nomfichier]
            Crée un dictionnaire mots\_répétés, dont les clés sont des mots d'une liste et les valeurs sont ces mots répétés trois fois. Voici comment fonctionne le code :

    words = ['apple', 'banana', 'cherry'] : Cette ligne initialise une liste nommée words avec trois valeurs de chaîne.
    repeated\_words = \{word : word * 3 for word in words\} : Cette ligne utilise une compréhension de dictionnaire pour créer le dictionnaire repeated\_words. Elle parcourt chaque mot de la liste words et crée une paire clé-valeur pour chaque mot.
        for word in words : Cette partie du code parcourt chaque mot de la liste des mots.
        'repeated\_words' = \{word : word * 3 for word in words\} : Pour chaque mot, le code prend le mot lui-même comme clé et utilise la multiplication de chaîne (mot * 3) pour répéter le mot trois fois, créant ainsi la valeur.
    print(words) : Cette ligne affiche la liste originale des mots sur la console.
    print(mots\_répétés) : Cette ligne affiche le dictionnaire mots\_répétés, qui contient les mots comme clés et les mots répétés trois fois comme valeurs.
        \end{solution}
        

        \question
        Créer un dictionnaire de chaînes avec les mots contenant "a" et leur longueur

Exemple de résultat

['apple', 'banana', 'cherry', 'date']

\{'apple' : 5, "banane" : 6, "date" : 4\}
        \par
        \begin{solution}
            \renewcommand{\nomfichier}{q685.py}
            \pythonfile{\chemincode \nomfichier}[][\nomfichier]
            Crée un dictionnaire a\_word\_lengths, dont les clés sont les mots d'une liste contenant la lettre "a" et les valeurs sont les longueurs de ces mots. Voici comment fonctionne le code :

    words = ['apple', 'banana', 'cherry', 'date'] : Cette ligne initialise une liste nommée words avec quatre valeurs de chaîne.
    a\_word\_lengths = \{word : len(word) for word in words if 'a' in word\} : Cette ligne utilise une compréhension de dictionnaire pour créer le dictionnaire a\_word\_lengths. Elle parcourt chaque mot de la liste des mots, vérifie si la lettre "a" est présente dans le mot et crée une paire clé-valeur pour les mots qui contiennent "a".
        pour mot dans mots : Cette partie du code parcourt chaque mot de la liste des mots.
        'a\_word\_lengths' = \{word : len(word) for word in words if 'a' in word\} : Pour chaque mot, il prend le mot lui-même comme clé et utilise la fonction len(word) pour calculer la longueur du mot, créant ainsi la valeur.
        if 'a' in word : Cette condition garantit que seuls les mots contenant la lettre 'a' sont inclus dans le dictionnaire.
    print(words) : Cette ligne affiche la liste originale des mots sur la console.
    print(a\_word\_lengths) : Cette ligne affiche le dictionnaire a\_word\_lengths, qui contient les mots (avec 'a') comme clés et leurs longueurs respectives comme valeurs.
        \end{solution}
        

        \question
        Créer un dictionnaire des nombres et de leurs carrés, mais uniquement pour les nombres supérieurs à 5

Exemple de résultat

[1, 2, 3, 4, 5, 6, 7, 8, 9, 10]

\{6 : 36, 7 : 49, 8 : 64, 9 : 81, 10 : 100\}
        \par
        \begin{solution}
            \renewcommand{\nomfichier}{q686.py}
            \pythonfile{\chemincode \nomfichier}[][\nomfichier]
            Crée un dictionnaire nommé squares\_greater\_than\_5, où les clés sont des nombres d'une liste qui sont supérieurs à 5, et les valeurs sont les carrés de ces nombres. Voici comment fonctionne le code :

    nombres = [1, 2, 3, 4, 5, 6, 7, 8, 9, 10] : Cette ligne initialise une liste nommée numbers avec dix valeurs entières.
    carrés\_plus\_que\_5 = \{x : x ** 2 for x in numbers if x > 5\} : Cette ligne utilise une compréhension de dictionnaire pour créer le dictionnaire squares\_greater\_than\_5. Elle parcourt chaque nombre de la liste des nombres et vérifie s'il est supérieur à 5. Si la condition est remplie, elle crée une paire clé-valeur dans le dictionnaire.
        pour x dans nombres : Cette partie du code parcourt chaque nombre de la liste des nombres.
        carrés\_plus\_que\_5 = \{x : x ** 2 for x in numbers if x > 5\} : Pour chaque nombre, le code prend le nombre lui-même comme clé et calcule le carré du nombre à l'aide de l'expression x ** 2, créant ainsi la valeur.
        if x > 5 : Cette condition garantit que seuls les nombres supérieurs à 5 sont inclus dans le dictionnaire.
    print(numbers) : Cette ligne imprime la liste originale des nombres sur la console.
    print(carrés\_plus\_grands\_que\_5) : Cette ligne affiche le dictionnaire squares\_greater\_than\_5, qui contient les nombres supérieurs à 5 comme clés et leurs carrés correspondants comme valeurs.
        \end{solution}
        

        \question
        Créer un dictionnaire de caractères et de leurs valeurs ASCII à partir d'une chaîne de caractères, à l'exclusion des caractères non alphabétiques.

Exemple de sortie

Bonjour123

\{'H' : 72, 'e' : 101, 'l' : 108, 'o' : 111\}
        \par
        \begin{solution}
            \renewcommand{\nomfichier}{q687.py}
            \pythonfile{\chemincode \nomfichier}[][\nomfichier]
            Crée un dictionnaire nommé ascii\_values pour faire correspondre les caractères alphabétiques de la variable texte à leurs valeurs ASCII. Voici comment fonctionne le code :

    text = "Hello123" : Cette ligne initialise la variable texte avec la chaîne "Hello123".
    ascii\\values = \{char : ord(char) for char in text if char.isalpha()\} : Cette ligne utilise une compréhension de dictionnaire pour créer le dictionnaire ascii\_values. Elle parcourt chaque caractère (char) de la chaîne de texte et vérifie si le caractère est alphabétique à l'aide de la méthode char.isalpha(). Si le caractère est alphabétique, il crée une paire clé-valeur dans le dictionnaire.
        pour char dans text : Cette partie du code parcourt chaque caractère de la chaîne de texte.
        ascii\_values = \{char : ord(char) for char in text if char.isalpha()\} : Pour chaque caractère alphabétique, il prend le caractère lui-même comme clé et récupère sa valeur ASCII à l'aide de la fonction ord(char), créant ainsi la valeur.
        if char.isalpha() : Cette condition garantit que seuls les caractères alphabétiques sont inclus dans le dictionnaire.
    print(text) : Cette ligne imprime la chaîne de texte originale sur la console.
    print(ascii\_values) : Cette ligne affiche le dictionnaire ascii\_values, qui contient des caractères alphabétiques comme clés et leurs valeurs ASCII correspondantes comme valeurs.
        \end{solution}
        

        \question
        Créer un dictionnaire de mots et de leurs formes en majuscules, en excluant les mots sans majuscules.

Exemple de résultat

['apple', 'Banana', 'cherry', 'Date']

\{'Banane' : 'BANANE', 'Date' : 'DATE'\}
        \par
        \begin{solution}
            \renewcommand{\nomfichier}{q688.py}
            \pythonfile{\chemincode \nomfichier}[][\nomfichier]
            Crée un dictionnaire nommé uppercase\_words qui associe les mots comportant au moins un caractère en majuscule à leur version en majuscule. Voici comment fonctionne le code :

    words = ['apple', 'Banana', 'cherry', 'Date'] : Cette ligne initialise la liste des mots avec une collection de mots, dont certains contiennent des caractères majuscules.
    uppercase\_words = \{word : word.upper() for word in words if any(char.isupper() for char in word)\} : Cette ligne utilise une compréhension de dictionnaire pour créer le dictionnaire uppercase\_words. Elle parcourt chaque mot de la liste des mots et vérifie s'il contient au moins un caractère majuscule (char.isupper()). Si le mot contient au moins une majuscule, il crée une paire clé-valeur dans le dictionnaire.
        pour mot dans mots : Cette partie du code parcourt chaque mot de la liste des mots.
        if any(char.isupper() for char in word) : Cette condition vérifie s'il existe au moins un caractère majuscule dans le mot actuel. La fonction any vérifie si l'un des caractères du mot est en majuscule.
        {word : word.upper() for word in words if any(char.isupper() for char in word)} : Si la condition est remplie, une paire clé-valeur est créée dans le dictionnaire. La clé est le mot original, et la valeur est la version en majuscules de ce mot obtenue à l'aide de word.upper().
    print(words) : Cette ligne affiche la liste des mots originaux sur la console.
    print(mots\_en\_majuscules) : Cette ligne affiche le dictionnaire uppercase\_words, qui contient des mots comportant au moins un caractère majuscule comme clés et leurs versions majuscules comme valeurs.
        \end{solution}
        

        \question
        Créer un dictionnaire de chaînes de mots contenant plus de 4 lettres

Exemple de résultat

Python est un langage de programmation puissant et polyvalent.

\{'Python' : 6, 'puissant' : 8, 'polyvalent' : 9, "programmation" : 11, "langage" : 8\}
        \par
        \begin{solution}
            \renewcommand{\nomfichier}{q689.py}
            \pythonfile{\chemincode \nomfichier}[][\nomfichier]
            Le code Python présente une phrase et crée un dictionnaire appelé mots\_longs. Ce dictionnaire associe les mots d'une longueur supérieure à 4 à leurs longueurs respectives. Voici comment fonctionne le code :

    phrase = "Python est un langage de programmation puissant et polyvalent" : Cette ligne initialise la variable phrase avec une chaîne contenant des mots.
    long\_words = \{word : len(word) for word in sentence.split() if len(word) > 4\} : Cette ligne utilise une compréhension de dictionnaire pour créer le dictionnaire long\_words. Elle divise la phrase en mots à l'aide de sentence.split(), puis parcourt chaque mot de la liste des mots.
        for word in sentence.split() : Cette partie du code parcourt chaque mot de la phrase.
        if len(word) > 4 : cette condition vérifie si la longueur du mot actuel est supérieure à 4 caractères.
        {mot : len(mot)} : Si la condition est remplie, une paire clé-valeur est créée dans le dictionnaire. La clé est le mot lui-même, et la valeur est la longueur du mot obtenue à l'aide de len(mot).
    print(phrase) : Cette ligne affiche la phrase originale sur la console.
    print(mots\_longs) : Cette ligne affiche le dictionnaire mots\_longs, qui contient des mots de longueur supérieure à 4 comme clés et leurs longueurs respectives comme valeurs.
        \end{solution}
        

        \question
        Créer un dictionnaire des mots et de leur première lettre en majuscule

Exemple de résultat

['apple', 'banana', 'cherry']

\{'apple' : 'Apple', 'banana' : 'Banana', 'cherry' : 'Cherry'\}
        \par
        \begin{solution}
            \renewcommand{\nomfichier}{q690.py}
            \pythonfile{\chemincode \nomfichier}[][\nomfichier]
            La liste des mots et crée un dictionnaire appelé capitalized\_first\_letter. Ce dictionnaire associe chaque mot à une version du mot dont la première lettre est en majuscule. Voici comment fonctionne le code :

    words = ['apple', 'banana', 'cherry'] : Cette ligne initialise la variable words avec une liste de mots.
    première\_lettre\_capitalisée = \{mot : mot[0].upper() + mot[1 :] pour mot dans mots\} : Cette ligne utilise une compréhension de dictionnaire pour créer le dictionnaire capitalized\_first\_letter. Elle parcourt chaque mot de la liste des mots.
        \{mot : mot[0].upper() + mot[1 :]\} : Cette partie du code crée des paires clé-valeur dans le dictionnaire. La clé est le mot original et la valeur est une version modifiée du mot dont la première lettre est en majuscule. Pour ce faire, on utilise word[0].upper() pour mettre la première lettre en majuscule, puis on la concatène avec le reste du mot obtenu à l'aide de word[1 :].
    print(words) : Cette ligne imprime la liste originale de mots, words, sur la console.
    print(première\_lettre\_capitalisée) : Cette ligne affiche le dictionnaire capitalized\_first\_letter, qui contient les mots originaux comme clés et leurs versions avec la première lettre en majuscule comme valeurs.
        \end{solution}
        

        \question
        Créer un dictionnaire de chaînes de caractères dont les voyelles ont été supprimées

Exemple de résultat

['apple', 'banana', 'cherry']

\{'apple' : 'ppl', 'banana' : 'bnn', 'cherry' : 'chrry'\}
        \par
        \begin{solution}
            \renewcommand{\nomfichier}{q691.py}
            \pythonfile{\chemincode \nomfichier}[][\nomfichier]
            La liste des mots et crée un dictionnaire nommé without\_vowels. Ce dictionnaire associe chaque mot à une version du mot dans laquelle toutes les voyelles (majuscules et minuscules) ont été supprimées. Voici comment fonctionne le code :

    words = ['apple', 'banana', 'cherry'] : Cette ligne initialise la variable words avec une liste de mots.
    without\_vowels = \{word : ''.join([char for char in word if char.lower() not in 'aeiou']) for word in words\} : Cette ligne utilise une compréhension de dictionnaire pour créer le dictionnaire sans\_voyelles. Elle parcourt chaque mot de la liste des mots.
        \{mot : ''.join([char for char in word if char.lower() not in 'aeiou'])\} : Cette partie du code crée des paires clé-valeur dans le dictionnaire. La clé est le mot original et la valeur est une version modifiée du mot. La modification est effectuée en parcourant chaque caractère char dans le mot et en vérifiant si char.lower() (la version minuscule du caractère) n'est pas dans la chaîne 'aeiou', qui représente les voyelles. Si char n'est pas une voyelle, il est inclus dans la version modifiée du mot. La fonction join est utilisée pour concaténer les caractères en une seule chaîne.
    print(words) : Cette ligne imprime la liste originale des mots, words, sur la console.
    print(sans\_voyelles) : Cette ligne affiche le dictionnaire without\_vowels, qui contient les mots originaux comme clés et leurs versions sans voyelles comme valeurs.
        \end{solution}
        

        \question
        Créer un dictionnaire des nombres avec leurs signes inversés

Exemple de résultat

[5, -10, 15, -20, 25]

\{5 : -5, -10 : 10, 15 : -15, -20 : 20, 25 : -25\}
        \par
        \begin{solution}
            \renewcommand{\nomfichier}{q692.py}
            \pythonfile{\chemincode \nomfichier}[][\nomfichier]
            La liste des nombres et crée un dictionnaire appelé signes\_inversés. Ce dictionnaire associe chaque nombre à sa négation (en changeant son signe de positif à négatif ou vice versa). Voici comment fonctionne le code :

    nombres = [5, -10, 15, -20, 25] : Cette ligne initialise la variable numbers avec une liste d'entiers, comprenant à la fois des nombres positifs et négatifs.
    signes\_inversés = \{x : -x for x in numbers\} : Cette ligne utilise une compréhension de dictionnaire pour créer le dictionnaire signes\_inversés. Elle parcourt chaque nombre x dans la liste des nombres.
        \{x : -x\} : Cette partie du code crée des paires clé-valeur dans le dictionnaire. La clé est le nombre original x, et la valeur est la négation de ce nombre, qui est calculée comme -x.
    print(nombres) : Cette ligne imprime la liste originale des nombres, numbers, sur la console.
    print(signes\_inversés) : Cette ligne affiche le dictionnaire signes\_inversés, qui contient les nombres originaux comme clés et leurs négations comme valeurs.
        \end{solution}
        

        \question
        Créer un dictionnaire de chaînes de caractères avec des mots inversés

Exemple de résultat

Python est amusant

\{0 : 'fun', 1 : 'is', 2 : 'Python'\}
        \par
        \begin{solution}
            \renewcommand{\nomfichier}{q693.py}
            \pythonfile{\chemincode \nomfichier}[][\nomfichier]
            Le code Python traite une phrase et crée un dictionnaire appelé mots\_inversés. Ce dictionnaire associe l'index (position) de chaque mot dans l'ordre inverse au mot lui-même. Voici comment fonctionne le code :

    phrase = "Python is fun" : Cette ligne initialise la variable phrase avec une chaîne contenant plusieurs mots.
    reversed\_words = \{i : word for i, word in enumerate(sentence.split()[::-1])\} : Cette ligne utilise une compréhension de dictionnaire pour créer le dictionnaire mots\_inversés. Voici ce qui se passe :
        sentence.split() : Cette partie du code divise la phrase en mots individuels et renvoie une liste de mots.
        [::-1] : Cette partie inverse l'ordre des mots dans la liste. Elle inverse donc l'ordre des mots dans la phrase.
        enumerate(...) : Cette fonction est utilisée pour parcourir la liste inversée des mots et renvoie à la fois l'index (i) et le mot lui-même.
        \{i : mot pour i, mot dans ...\} : Cette partie du code crée des paires clé-valeur dans le dictionnaire. Les clés (i) sont les indices des mots, et les valeurs (mot) sont les mots individuels de la liste inversée.
    print(phrase) : Cette ligne imprime la phrase originale, sentence, sur la console.
    print(mots\_inversés) : Cette ligne affiche le dictionnaire mots\_inversés, qui contient les indices des mots comme clés (dans l'ordre inverse) et les mots correspondants comme valeurs.
        \end{solution}
        

        \question
        Créer un dictionnaire des nombres avec leurs diviseurs

Exemple de résultat

[1, 2, 3, 4, 5]

\{1 : [1], 2 : [1, 2], 3 : [1, 3], 4 : [1, 2, 4], 5 : [1, 5]\}
        \par
        \begin{solution}
            \renewcommand{\nomfichier}{q694.py}
            \pythonfile{\chemincode \nomfichier}[][\nomfichier]
            Crée un dictionnaire nommé diviseurs qui associe chaque nombre de la liste des nombres à une liste de ses diviseurs. Voici comment fonctionne le code :

    nombres = [1, 2, 3, 4, 5] : Cette ligne initialise la liste des nombres avec une séquence d'entiers.
    diviseurs = \{x : [i for i in range(1, x + 1) if x \% i == 0] for x in numbers\} : Cette ligne utilise une compréhension de dictionnaire pour créer le dictionnaire des diviseurs. Voici ce qui se passe :
        for x in numbers : Cette partie du code itère sur chaque nombre x dans la liste des nombres.
        \{x : ...\} : Cette partie du code crée des paires clé-valeur dans le dictionnaire. La clé est le nombre x et la valeur est le résultat de la compréhension de la liste intérieure.
        [i for i in range(1, x + 1) if x \% i == 0] : Il s'agit de la compréhension de liste qui génère une liste de diviseurs pour chaque nombre x. Elle parcourt les nombres de 1 à x (inclus) et vérifie si x est divisible par i. Si c'est le cas, i est inclus dans la liste des diviseurs.
    print(nombres) : Cette ligne imprime la liste des nombres sur la console.
    print(diviseurs) : Cette ligne affiche le dictionnaire des diviseurs, qui contient des nombres comme clés et des listes de diviseurs comme valeurs.
        \end{solution}
        

        \question
        Créer un dictionnaire de caractères et leur nombre, à l'exclusion des caractères d'espacement, à partir d'une chaîne de caractères.

Exemple de sortie

bonjour le monde

\{'o' : 2, 'h' : 1, 'r' : 1, 'w' : 1, 'e' : 1, 'l' : 3, 'd' : 1\}
        \par
        \begin{solution}
            \renewcommand{\nomfichier}{q695.py}
            \pythonfile{\chemincode \nomfichier}[][\nomfichier]
            Il compte les occurrences de caractères non spatiaux dans le texte donné et crée un dictionnaire qui associe chaque caractère à son nombre. Voici comment fonctionne le code :

    text = "hello world" : Cette ligne initialise la variable texte avec le texte d'entrée.
    char\_counts = \{char : text.count(char) for char in set(text) if not char.isspace()\} : Cette ligne utilise une compréhension de dictionnaire pour créer le dictionnaire char\_counts. Voici ce qui se passe :
        \{char : ...\} : Cette partie du code crée des paires clé-valeur dans le dictionnaire. La clé est un caractère char, et la valeur est le résultat de l'expression text.count(char), qui compte les occurrences de char dans le texte.
        for char in set(text) : Cette partie itère sur chaque caractère unique du texte. L'expression set(text) est utilisée pour obtenir un ensemble de caractères uniques, en éliminant les doublons.
        if not char.isspace() : Cette partie vérifie si le caractère char n'est pas un caractère d'espacement. Elle permet de s'assurer que seuls les caractères non spatiaux sont comptés.
    print(text) : Cette ligne imprime le texte original sur la console.
    print(char\_counts) : Cette ligne affiche le dictionnaire char\_counts, qui contient des caractères non spatiaux comme clés et leur nombre comme valeurs.
        \end{solution}
        

        \question
        Correspondance entre les lettres minuscules et leurs valeurs ASCII

Exemple de sortie

\{'a' : 97, 'b' : 98, 'c' : 99, 'd' : 100, 'e' : 101, 'f' : 102, "g" : 103, "h" : 104, 'i' : 105, 'j' : 106, 'k' : 107, "l" : 108, "m" : 109, "n" : 110, "o" : 111, "p" : 112, 'q' : 113, 'r' : 114, 's' : 115, "t" : 116, "u" : 117, 'v' : 118, 'w' : 119, 'x' : 120, 'y' : 121, 'z' : 122\}
        \par
        \begin{solution}
            \renewcommand{\nomfichier}{q696.py}
            \pythonfile{\chemincode \nomfichier}[][\nomfichier]
            Crée un dictionnaire appelé ascii\_mapping qui associe chaque lettre minuscule de l'alphabet anglais à la valeur ASCII correspondante. Voici ce que fait le code :

    ascii\_mapping = \{char : ord(char) for char in 'abcdefghijklmnopqrstuvwxyz'\} : Cette ligne crée le dictionnaire ascii\_mapping en utilisant une compréhension de dictionnaire. Voici comment cela fonctionne :
        \{char : ord(char) for char in 'abcdefghijklmnopqrstuvwxyz'\} est la compréhension du dictionnaire. Il parcourt chaque caractère (char) de la chaîne 'abcdefghijklmnopqrstuvwxyz'.
        Pour chaque caractère, il crée une paire clé-valeur dans le dictionnaire. La clé (char) est le caractère lui-même, et la valeur (ord(char)) est la valeur ASCII de ce caractère obtenue à l'aide de la fonction ord.
        \end{solution}
        

        \question
        Nombres et leur factorielle de 1 à 10

Exemple de résultat

\{1 : 1, 2 : 2, 3 : 6, 4 : 24, 5 : 120, 6 : 720, 7 : 5040, 8 : 40320, 9 : 362880, 10 : 3628800\}
        \par
        \begin{solution}
            \renewcommand{\nomfichier}{q697.py}
            \pythonfile{\chemincode \nomfichier}[][\nomfichier]
            Le code définit une fonction Python factorial(n) pour calculer la factorielle d'un nombre n à l'aide de la récursivité, puis utilise cette fonction pour créer un dictionnaire appelé factorials. Ce dictionnaire associe les nombres de 1 à 10 à leurs valeurs factorielles respectives. Voici une explication pas à pas du code :

    def factorial(n) : Cette ligne définit une fonction nommée factorial qui prend un entier n en entrée et calcule la factorielle de n en utilisant la récursivité.
        Le cas de base est défini comme suit : if n == 0 :. Lorsque n est 0, la fonction renvoie 1 car 0 ! (lire "factorielle zéro") est défini comme 1.
        Dans le cas récursif, la fonction calcule la factorielle de n comme n * factorielle(n - 1). Elle s'appelle récursivement avec une valeur plus petite jusqu'à ce qu'elle atteigne le cas de base.
    factorials = \{x : factorial(x) for x in range(1, 11)\} : Cette ligne crée le dictionnaire factoriel en utilisant une compréhension de dictionnaire. Voici comment cela fonctionne :
        {x : factorial(x) for x in range(1, 11)} est la compréhension du dictionnaire. Elle parcourt la plage de nombres de 1 à 10 (inclus).
        Pour chaque nombre x, il crée une paire clé-valeur dans le dictionnaire. La clé (x) est le nombre lui-même, et la valeur (factorielle(x)) est calculée en appelant la fonction factorielle avec ce nombre en entrée.
    print(factorials) : Cette ligne affiche le dictionnaire factoriel sur la console.
        \end{solution}
        

        \question
        Les nombres et leur représentation binaire de 1 à 10

Exemple de sortie

\{1 : '1', 2 : '10', 3 : '11', 4 : '100', 5 : '101', 6 : '110', 7 : '111', 8 : '1000', 9 : '1001', 10 : '1010'\}
        \par
        \begin{solution}
            \renewcommand{\nomfichier}{q698.py}
            \pythonfile{\chemincode \nomfichier}[][\nomfichier]
            Crée un dictionnaire appelé binary\_rep en utilisant une compréhension de dictionnaire en Python. Ce dictionnaire associe les nombres de 1 à 10 à leurs représentations binaires respectives sous forme de chaînes de caractères. Voici une explication pas à pas du code :

    binary\_rep = \{x : bin(x)[2 :] for x in range(1, 11)\} : Cette ligne crée le dictionnaire binary\_rep à l'aide d'une compréhension de dictionnaire. Voici comment cela fonctionne :
        \{x : bin(x)[2 :] for x in range(1, 11)\} est la compréhension du dictionnaire. Elle parcourt la plage de nombres de 1 à 10 (inclus).
        Pour chaque nombre x, il crée une paire clé-valeur dans le dictionnaire. La clé (x) est le nombre lui-même, et la valeur (bin(x)[2 :]) est calculée en convertissant x en sa représentation binaire à l'aide de la fonction bin, puis en supprimant les deux premiers caractères (qui sont '0b') de la chaîne binaire pour obtenir la représentation binaire sous la forme d'une chaîne.
    print(binary\_rep) : Cette ligne affiche le dictionnaire binary\_rep sur la console.
        \end{solution}
        

        \question
        Paires d'éléments distincts et leur différence absolue à partir de deux listes

Exemple de résultat

[3, 6, 9]

[5, 10, 15]

\{(3, 5) : 2, (3, 10) : 7, (3, 15) : 12, (6, 5) : 1, (6, 10) : 4, (6, 15) : 9, (9, 5) : 4, (9, 10) : 1, (9, 15) : 6\}
        \par
        \begin{solution}
            \renewcommand{\nomfichier}{q699.py}
            \pythonfile{\chemincode \nomfichier}[][\nomfichier]
            Crée un dictionnaire nommé abs\_diff\_dict. Ce dictionnaire associe des paires d'éléments de list1 et list2 à leurs différences absolues. Voici comment fonctionne le code :

    list1 = [3, 6, 9] et list2 = [5, 10, 15] : Ces lignes initialisent deux listes, list1 et list2, avec des valeurs entières.
    abs\_diff\_dict = \{(x, y) : abs(x - y) for x in list1 for y in list2\} : Cette ligne utilise une compréhension de dictionnaire pour créer le dictionnaire abs\_diff\_dict. Voici ce qui se passe :
        (x, y) est utilisé comme un tuple pour représenter des paires d'éléments, où x est un élément de la liste 1 et y est un élément de la liste 2.
        abs(x - y) calcule la différence absolue entre x et y pour chaque paire.
        La compréhension du dictionnaire passe en revue toutes les paires possibles d'éléments de list1 et list2 et calcule la différence absolue pour chaque paire, en utilisant le tuple (x, y) comme clé et la différence absolue comme valeur.
    print(list1) : Cette ligne affiche le contenu de la liste 1 sur la console.
    print(list2) : Cette ligne affiche le contenu de la liste 2 sur la console.
    print(abs\_diff\_dict) : Cette ligne affiche le dictionnaire abs\_diff\_dict, qui contient des paires d'éléments comme clés et leurs différences absolues comme valeurs.
        \end{solution}
        

        \question
        Paires d'éléments distincts et leur position de caractère additionnés à partir de deux listes

Exemple de résultat

['abc', 'def', 'ghi']

['jkl', 'mno', 'pqr']

\{('abc', 'jkl') : 615, ('abc', 'mno') : 624, ('abc', 'pqr') : 633, ('def', 'jkl') : 624, ('def', 'mno') : 633, ("def", "pqr") : 642, ('ghi', 'jkl') : 633, ('ghi', 'mno') : 642, ('ghi', 'pqr') : 651\}
        \par
        \begin{solution}
            \renewcommand{\nomfichier}{q700.py}
            \pythonfile{\chemincode \nomfichier}[][\nomfichier]
            Crée un dictionnaire appelé char\_pos\_sum\_dict à l'aide d'une compréhension de dictionnaire en Python. Ce dictionnaire associe des paires de chaînes de caractères provenant de deux listes, list1 et list2, à la somme des valeurs ordinales (valeurs ASCII) des caractères de chaque paire. Voici une explication pas à pas du code :

    list1 = ["abc", "def", "ghi"] : Cette ligne définit une liste appelée list1 contenant trois chaînes de caractères.
    list2 = ["jkl", "mno", "pqr"] : Cette ligne définit une autre liste appelée list2 contenant trois chaînes de caractères.
    char\_pos\_sum\_dict = \{(x, y) : sum(ord(char) for char in x) + sum(ord(char) for char in y) for x in list1 for y in list2\} : Cette ligne crée le dictionnaire char\_pos\_sum\_dict à l'aide d'une compréhension de dictionnaire imbriqué. Voici comment cela fonctionne :
        \{(x, y) : sum(ord(char) for char in x) + sum(ord(char) for char in y) for x in list1 for y in list2\} est la compréhension du dictionnaire. Il itère sur des paires de chaînes (x, y) où x provient de list1 et y de list2.
        Pour chaque paire de chaînes, il crée une paire clé-valeur dans le dictionnaire. La clé x, y est un tuple contenant les deux chaînes. La valeur est calculée en additionnant les valeurs ordinales (valeurs ASCII) de tous les caractères de la première chaîne x et de tous les caractères de la deuxième chaîne y.
    print(list1) : Cette ligne imprime le contenu de la liste 1 sur la console.
    print(list2) : Cette ligne affiche le contenu de la liste 2 sur la console.
    print(char\_pos\_sum\_dict) : Cette ligne affiche le dictionnaire char\_pos\_sum\_dict sur la console.
        \end{solution}
        

        \question
        Mots distincts et leur longueur, à l'exclusion des mots de longueur impaire, dans une phrase

Exemple de sortie

Bonjour, comment allez-vous ?

\{'Hello,' : 6, 'you?' : 4\}
        \par
        \begin{solution}
            \renewcommand{\nomfichier}{q701.py}
            \pythonfile{\chemincode \nomfichier}[][\nomfichier]
            Le code traite une phrase et crée un dictionnaire nommé distinct\_word\_length\_no\_odd à l'aide d'une compréhension de dictionnaire en Python. Ce dictionnaire associe les mots distincts de la phrase à leurs longueurs respectives, mais uniquement pour les mots de longueur paire. Voici une explication pas à pas du code :

    sentence = "Hello, how are you ?": Cette ligne définit une variable de type chaîne appelée sentence contenant la phrase "Hello, how are you ?"
    distinct\_word\_length\_no\_odd = \{word : len(word) for word in set(sentence.split()) if len(word) \% 2 == 0\} : Cette ligne crée le dictionnaire distinct\_word\_length\_no\_odd en utilisant une compréhension de dictionnaire. Voici comment cela fonctionne :
        \{word : len(word) for word in set(sentence.split()) if len(word) \% 2 == 0\} est la compréhension du dictionnaire. Il effectue les étapes suivantes :
        sentence.split() divise la phrase en une liste de mots. Dans ce cas, elle inclura des mots comme "Hello," (avec une virgule) et "you ?" (avec un point d'interrogation). (avec un point d'interrogation).
        set(sentence.split()) crée un ensemble de mots distincts en supprimant les doublons et la ponctuation. Il contiendra donc des mots comme "Hello," (sans la virgule) et "you" (sans le point d'interrogation).
        for word in set(sentence.split()) itère sur chaque mot distinct de l'ensemble.
        len(word) \% 2 == 0 vérifie si la longueur du mot est paire.
        Si la longueur d'un mot est paire, il crée une paire clé-valeur dans le dictionnaire. La clé (mot) est le mot lui-même, et la valeur (len(mot)) est la longueur du mot.
    print(phrase) : Cette ligne imprime la phrase originale, qui est "Hello, how are you ?", sur la console.
    print(distinct\_word\_length\_no\_odd) : Cette ligne affiche le dictionnaire distinct\_mot\_longueur\_no\_odd sur la console.
        \end{solution}
        

        \question
        Mots distincts et leur longueur, à l'exclusion des mots de longueur paire, dans une phrase

Exemple de sortie

Bonjour, comment allez-vous ?

\{'how' : 3, 'are' : 3\}
        \par
        \begin{solution}
            \renewcommand{\nomfichier}{q702.py}
            \pythonfile{\chemincode \nomfichier}[][\nomfichier]
            Le code traite une phrase et crée un dictionnaire nommé distinct\_word\_length\_no\_even à l'aide d'une compréhension de dictionnaire en Python. Ce dictionnaire associe les mots distincts de la phrase à leurs longueurs respectives, mais uniquement pour les mots de longueur impaire. Voici une explication pas à pas du code :

    sentence = "Hello, how are you ?": Cette ligne définit une variable de type chaîne appelée sentence contenant la phrase "Hello, how are you ?"
    distinct\_word\_length\_no\_even = \{word : len(word) for word in set(sentence.split()) if len(word) \% 2 != 0\} : Cette ligne crée le dictionnaire distinct\_mot\_longueur\_même à l'aide d'une compréhension de dictionnaire. Voici comment cela fonctionne :
        \{word : len(word) for word in set(sentence.split()) if len(word) \% 2 != 0\} est la compréhension du dictionnaire. Il effectue les étapes suivantes :
        sentence.split() divise la phrase en une liste de mots. Dans ce cas, elle inclura des mots comme "Hello," (avec une virgule) et "you ?" (avec un point d'interrogation). (avec un point d'interrogation).
        set(sentence.split()) crée un ensemble de mots distincts en supprimant les doublons et la ponctuation. Il contiendra donc des mots comme "Hello," (sans la virgule) et "you" (sans le point d'interrogation).
        for word in set(sentence.split()) itère sur chaque mot distinct de l'ensemble.
        len(word) \% 2 != 0 vérifie si la longueur du mot est impaire.
        Si la longueur d'un mot est impaire, cela crée une paire clé-valeur dans le dictionnaire. La clé (mot) est le mot lui-même, et la valeur (len(mot)) est la longueur du mot.
    print(phrase) : Cette ligne imprime la phrase originale, qui est "Hello, how are you ?", sur la console.
    print(distinct\_word\_length\_no\_even) : Cette ligne affiche le dictionnaire distinct\_mot\_longueur\_non\_pair sur la console.
        \end{solution}
        

        \question
        Caractères et nombre d'occurrences, hors ponctuation, dans une phrase

Exemple de sortie

Bonjour, comment allez-vous ?

\{'w' : 1, 'r' : 1, 'u' : 1, 'o' : 3, 'e' : 2, 'H' : 1, 'l' : 2, 'h' : 1, ' ' : 3, 'y' : 1, 'a' : 1\}
        \par
        \begin{solution}
            \renewcommand{\nomfichier}{q703.py}
            \pythonfile{\chemincode \nomfichier}[][\nomfichier]
            Le code traite une phrase et crée un dictionnaire nommé char\_occurrence\_no\_punct à l'aide d'une compréhension de dictionnaire en Python. Ce dictionnaire associe les caractères présents dans la phrase à leur nombre respectif, à l'exclusion des caractères de ponctuation. Voici une explication pas à pas du code :

    import string : Cette ligne importe le module string, qui fournit une chaîne contenant tous les caractères de ponctuation ASCII.
    sentence = "Hello, how are you ?": Cette ligne définit une variable de type chaîne de caractères appelée sentence contenant la phrase "Hello, how are you ?"
    char\_occurrence\_no\_punct = \{char : sentence.count(char) for char in set(sentence) if char not in string.punctuation\} : Cette ligne crée le dictionnaire char\_occurrence\_no\_punct à l'aide d'une compréhension de dictionnaire. Voici comment cela fonctionne :
        \{char : sentence.count(char) for char in set(sentence) if char not in string.punctuation\} est la compréhension du dictionnaire. Il exécute les étapes suivantes :
        set(sentence) crée un ensemble de caractères distincts à partir de la phrase.
        for char in set(sentence) itère sur chaque caractère distinct de l'ensemble.
        char not in string.punctuation vérifie si le caractère n'est pas dans la chaîne string.punctuation, qui contient tous les caractères de ponctuation.
        Si un caractère n'est pas un caractère de ponctuation, il crée une paire clé-valeur dans le dictionnaire. La clé (char) est le caractère lui-même, et la valeur (sentence.count(char)) est le nombre de fois où ce caractère apparaît dans la phrase.
    print(phrase) : Cette ligne imprime la phrase originale, qui est "Hello, how are you ?", sur la console.
    print(char\_occurrence\_no\_punct) : Cette ligne affiche le dictionnaire char\_occurrence\_no\_punct sur la console.
        \end{solution}
        

        \question
        Mots distincts et leur longueur, à l'exclusion de ceux dont la longueur n'est pas de Fibonacci, dans une phrase

Exemple de sortie

Bonjour, comment allez-vous ?

\{'are' : 3, 'how' : 3\}
        \par
        \begin{solution}
            \renewcommand{\nomfichier}{q704.py}
            \pythonfile{\chemincode \nomfichier}[][\nomfichier]
            Le code définit une fonction is\_fibonacci\_length(n) pour vérifier si un nombre donné est un nombre de Fibonacci (en vérifiant si sa longueur fait partie de la séquence de Fibonacci). Il utilise ensuite cette fonction pour créer un dictionnaire nommé distinct\_word\_length\_fibonacci à l'aide d'une compréhension de dictionnaire en Python. Ce dictionnaire associe les mots distincts de la phrase à leurs longueurs respectives, mais uniquement pour les mots dont les longueurs sont des nombres de Fibonacci. Voici une explication pas à pas du code :

    def is\_fibonacci\_length(n) : Cette ligne définit une fonction nommée is\_fibonacci\_length qui prend un entier n en entrée et vérifie si n est un nombre de Fibonacci en comparant sa longueur à la séquence de Fibonacci.
        Dans la fonction, une liste fib est initialisée avec les deux premiers nombres de Fibonacci (0 et 1).
        Une boucle while vérifie si le dernier nombre de la liste fib (fib[-1]) est inférieur ou égal à n. Elle continue à calculer les nombres de Fibonacci jusqu'à ce que le dernier soit inférieur ou égal à n.
        Si fib[-1] est égal à n, la fonction renvoie True car n est un nombre de Fibonacci en raison de sa longueur.
        Si la boucle se termine sans trouver de nombre de Fibonacci de la longueur donnée, la fonction renvoie False.
    sentence = "Hello, how are you ?": Cette ligne définit une variable de type chaîne de caractères appelée sentence contenant la phrase "Hello, how are you ?"
    distinct\_word\_length\_fibonacci = {word : len(word) for word in set(sentence.split()) if is\_fibonacci\_length(len(word))} : Cette ligne crée le dictionnaire distinct\_mot\_longueur\_fibonacci à l'aide d'une compréhension de dictionnaire. Voici comment cela fonctionne :
        \{word : len(word) for word in set(sentence.split()) if is\_fibonacci\_length(len(word)\} est la compréhension du dictionnaire. Il effectue les étapes suivantes :
        sentence.split() divise la phrase en une liste de mots.
        set(sentence.split()) crée un ensemble de mots distincts en supprimant les doublons.
        for word in set(sentence.split()) itère sur chaque mot distinct de l'ensemble.
        is\_fibonacci\_length(len(word)) vérifie si la longueur du mot est un nombre de Fibonacci à l'aide de la fonction is\_fibonacci\_length.
        Si la longueur d'un mot est un nombre de Fibonacci, une paire clé-valeur est créée dans le dictionnaire. La clé (mot) est le mot lui-même, et la valeur (len(mot)) est la longueur du mot.
    print(phrase) : Cette ligne imprime la phrase originale, qui est "Hello, how are you ?", sur la console.
    print(longueur\_de\_mot\_distinct\_fibonacci) : Cette ligne affiche le dictionnaire distinct\_mot\_longueur\_fibonacci sur la console.
        \end{solution}
        

        \question
        Correspondance entre les mots et la fréquence à laquelle ils contiennent la lettre "e".

Exemple de résultat

Bonjour, comment allez-vous ?

\{'Hello,' : 1, 'comment' : 0, 'are' : 1, 'you ? 0\}
        \par
        \begin{solution}
            \renewcommand{\nomfichier}{q705.py}
            \pythonfile{\chemincode \nomfichier}[][\nomfichier]
            Le code traite une phrase et crée un dictionnaire nommé word\_contains\_e à l'aide d'une compréhension de dictionnaire en Python. Ce dictionnaire associe les mots de la phrase au nombre de lettres "e" dans chaque mot. Voici une explication pas à pas du code :

    sentence = "Hello, how are you ?": Cette ligne définit une variable de type chaîne appelée sentence contenant la phrase "Hello, how are you ?"
    word\_contains\_e = \{word : word.count('e') for word in sentence.split()\} : Cette ligne crée le dictionnaire word\_contains\_e à l'aide d'une compréhension de dictionnaire. Voici comment cela fonctionne :
        \{word : word.count('e') for word in sentence.split()\} est la compréhension du dictionnaire. Elle effectue les étapes suivantes :
        sentence.split() divise la phrase en une liste de mots. Dans ce cas, elle inclura des mots comme "Hello," (avec une virgule) et "you ?" (avec un point d'interrogation). (avec un point d'interrogation).
        for word in sentence.split() parcourt chaque mot de la liste de mots.
        word.count('e') compte les occurrences de la lettre "e" dans le mot.
        Il crée une paire clé-valeur dans le dictionnaire, où la clé (mot) est le mot lui-même, et la valeur (mot.count('e')) est le nombre d'occurrences de la lettre 'e' dans le mot.
    print(phrase) : Cette ligne imprime la phrase originale, qui est "Hello, how are you ?", sur la console.
    print(mot\_contient\_e) : Cette ligne affiche le dictionnaire word\_contains\_e sur la console.
        \end{solution}
        

        \question
        Correspondance entre les caractères et leur position dans l'alphabet

Exemple de sortie

\{'a' : 1, 'b' : 2, 'c' : 3, 'd' : 4, 'e' : 5, 'f' : 6, 'g' : 7, 'h' : 8, 'i' : 9, 'j' : 10, 'k' : 11, "l" : 12, "m" : 13, "n" : 14, "o" : 15, "p" : 16, 'q' : 17, 'r' : 18, 's' : 19, "t" : 20, "u" : 21, "v" : 22, "w" : 23, 'x' : 24, 'y' : 25, 'z' : 26\}
        \par
        \begin{solution}
            \renewcommand{\nomfichier}{q706.py}
            \pythonfile{\chemincode \nomfichier}[][\nomfichier]
            Crée un dictionnaire nommé char\_alphabet\_position à l'aide d'une compréhension de dictionnaire en Python. Ce dictionnaire associe les lettres minuscules de l'alphabet anglais à leurs positions respectives dans l'alphabet, en commençant par 1. Voici une explication du code étape par étape :

    import string : Cette ligne importe le module string, qui fournit une constante string.ascii\_lowercase contenant toutes les lettres minuscules de l'alphabet anglais.
    char\_alphabet\_position = {char : string.ascii\_lowercase.index(char) + 1 for char in string.ascii\_lowercase} : Cette ligne crée le dictionnaire char\_alphabet\_position à l'aide d'une compréhension de dictionnaire. Voici comment cela fonctionne :
        \{char : string.ascii\_lowercase.index(char) + 1 for char in string.ascii\_lowercase\} est la compréhension du dictionnaire. Il effectue les étapes suivantes :
        string.ascii\_lowercase est la chaîne contenant toutes les lettres minuscules de l'alphabet anglais : "abcdefghijklmnopqrstuvwxyz".
        for char in string.ascii\_lowercase itère sur chaque lettre minuscule de l'alphabet.
        string.ascii\_lowercase.index(char) trouve l'indice (position) du caractère char dans la chaîne de caractères string.ascii\_lowercase. Cet indice est basé sur zéro.
        + 1 est ajouté à l'index pour le faire passer d'un index basé sur zéro à une position basée sur un, car les positions dans l'alphabet commencent à 1.
        Cela crée une paire clé-valeur dans le dictionnaire, où la clé (char) est la lettre minuscule, et la valeur (string.ascii\_lowercase.index(char) + 1) est sa position dans l'alphabet.
    print(char\_alphabet\_position) : Cette ligne affiche le dictionnaire char\_alphabet\_position sur la console.
        \end{solution}
        

        \question
        Mise en correspondance des mots avec leur statut de palindrome dans une phrase

Exemple de sortie

madame voit la voiture de course

\{'madam' : True, 'sees' : True, 'the' : False, 'racecar' : True\}
        \par
        \begin{solution}
            \renewcommand{\nomfichier}{q707.py}
            \pythonfile{\chemincode \nomfichier}[][\nomfichier]
            Le code traite une phrase et crée un dictionnaire nommé word\_is\_palindrome à l'aide d'une compréhension de dictionnaire en Python. Ce dictionnaire associe les mots de la phrase à une valeur booléenne indiquant si chaque mot est un palindrome. Voici une explication pas à pas du code :

    sentence = "madame voit la voiture de course" : Cette ligne définit une variable de type chaîne appelée phrase contenant la phrase "madame voit la voiture de course".
    word\_is\_palindrome = \{word : word == word[::-1] for word in sentence.split()\} : Cette ligne crée le dictionnaire word\_is\_palindrome en utilisant la compréhension du dictionnaire. Voici comment cela fonctionne :
        \{word : word == word[::-1] for word in sentence.split()\} est la compréhension du dictionnaire. Elle effectue les étapes suivantes :
        sentence.split() divise la phrase en une liste de mots. Dans ce cas, la phrase contient des mots comme "madam", "sees", "the" et "racecar".
        for word in sentence.split() itère sur chaque mot de la liste de mots.
        mot == mot[::-1] vérifie si le mot est un palindrome. Pour ce faire, il compare le mot lui-même (word) à son inverse (word[::-1]).
        Si un mot est un palindrome, il crée une paire clé-valeur dans le dictionnaire. La clé (mot) est le mot lui-même, et la valeur (mot == mot[::-1]) est True, indiquant qu'il s'agit d'un palindrome.
        Si un mot n'est pas un palindrome, la valeur est False.
    print(phrase) : Cette ligne imprime la phrase originale, qui est "madame voit la voiture de course", sur la console.
    print(mot\_est\_palindrome) : Cette ligne affiche le dictionnaire word\_is\_palindrome sur la console.
        \end{solution}
        

        \question
        Correspondance entre les nombres et leurs représentations binaires et hexadécimales

Exemple de sortie

\{1 : ('1', '1'), 2 : ('10', '2'), 3 : ('11', '3'), 4 : ('100', '4'), 5 : ('101', '5'), 6 : ('110', '6'), 7 : ('111', '7'), 8 : ('1000', '8'), 9 : ('1001', '9'), 10 : ('1010', 'a')\}
        \par
        \begin{solution}
            \renewcommand{\nomfichier}{q708.py}
            \pythonfile{\chemincode \nomfichier}[][\nomfichier]
            Crée un dictionnaire nommé number\_binary\_hex en utilisant une compréhension de dictionnaire en Python. Ce dictionnaire associe les nombres de 1 à 10 à leurs représentations binaires et hexadécimales. Voici une explication pas à pas du code :

    number\_binary\_hex = \{num : (bin(num)[2 :], hex(num)[2 :]) for num in range(1, 11)\} : Cette ligne crée le dictionnaire number\_binary\_hex à l'aide d'une compréhension de dictionnaire. Voici comment cela fonctionne :
        \{num : (bin(num)[2 :], hex(num)[2 :]) for num in range(1, 11)\} est la compréhension du dictionnaire. Il effectue les étapes suivantes :
        for num in range(1, 11) itère sur les nombres de 1 à 10 (inclus).
        (bin(num)[2 :], hex(num)[2 :]) est un tuple contenant deux éléments :
            bin(num)[2 :] convertit le nombre num en sa représentation binaire à l'aide de la fonction bin, puis coupe les deux premiers caractères (qui représentent "0b" au début de la chaîne binaire).
            hex(num)[2 :] convertit le nombre num en sa représentation hexadécimale à l'aide de la fonction hex, puis coupe les deux premiers caractères (qui représentent "0x" au début de la chaîne hexadécimale).
        Il crée une paire clé-valeur dans le dictionnaire, où la clé (num) est le nombre lui-même, et la valeur est le tuple contenant les représentations binaire et hexadécimale.
    print(nombre\_binaire\_hex) : Cette ligne imprime le dictionnaire number\_binary\_hex sur la console.
        \end{solution}
        

        \question
        Mettre en correspondance les mots et leur inverse dans une phrase

Exemple de sortie

Bonjour, comment allez-vous ?

\{'Hello,' : ',olleH', 'how' : 'woh', 'are' : 'era', 'you?' : '?uoy'\}
        \par
        \begin{solution}
            \renewcommand{\nomfichier}{q709.py}
            \pythonfile{\chemincode \nomfichier}[][\nomfichier]
            Le code traite une phrase et crée un dictionnaire nommé word\_reverse\_mapping à l'aide d'une compréhension de dictionnaire en Python. Ce dictionnaire associe les mots de la phrase à leurs versions inversées respectives. Voici une explication pas à pas du code :

    sentence = "Hello, how are you ?": Cette ligne définit une variable de chaîne appelée sentence contenant la phrase "Hello, how are you ?"
    word\_reverse\_mapping = \{word : word[::-1] for word in sentence.split()\} : Cette ligne crée le dictionnaire word\_reverse\_mapping à l'aide d'une compréhension de dictionnaire. Voici comment cela fonctionne :
        {word : word[::-1] for word in sentence.split()} est la compréhension du dictionnaire. Elle effectue les étapes suivantes :
        sentence.split() divise la phrase en une liste de mots. Dans ce cas, elle inclura des mots comme "Hello," (avec une virgule) et "you ?" (avec un point d'interrogation). (avec un point d'interrogation).
        for word in sentence.split() parcourt chaque mot de la liste de mots.
        mot[::-1] inverse l'ordre des caractères dans le mot. Pour ce faire, on utilise la notation de découpage de Python, où [::-1] est utilisé pour inverser une chaîne de caractères.
        Cela crée une paire clé-valeur dans le dictionnaire, où la clé (word) est le mot lui-même, et la valeur (word[::-1]) est le mot inversé.
    print(phrase) : Cette ligne imprime la phrase originale, qui est "Hello, how are you ?", sur la console.
    print(word\_reverse\_mapping) : Cette ligne affiche le dictionnaire word\_reverse\_mapping sur la console.
        \end{solution}
        

	
\end{questions}
%%*********************************************
%\section{Site 8 Class}
%\href{https://www.tutorjoes.in/python_programming_tutorial/class_objects_programs_in_python#google_vignette}{Dictionnary de comprehension}\par
%
%\begin{questions}
%	
%	
        \question
        DEPART de la classe
Ecrire un programme python pour créer une classe d'étudiants
        \par
        \begin{solution}
            \renewcommand{\nomfichier}{q710.py}
            \pythonfile{\chemincode \nomfichier}[][\nomfichier]
            Ce programme Python définit une classe Student avec des attributs et des méthodes permettant de définir et d'afficher les données des étudiants. Voici une explication du programme :

    classe Étudiant : Il s'agit de la classe principale représentant un étudiant. Elle possède les attributs et méthodes suivants :
        Attributs de la classe :
            id : Permet de stocker l'identifiant de l'étudiant (initialisé à 0 par défaut).
            name : Stocke le nom de l'étudiant (initialisé par défaut à une chaîne vide).
            gender (sexe) : Indique le sexe de l'étudiant (initialisé par défaut comme une chaîne vide).
            total : Stocke les notes totales de l'étudiant (initialisé par défaut comme une chaîne vide).
            per : Stocke le pourcentage de l'étudiant (initialisé à 0 par défaut).
        setData(self, id, name, gender, total, per) : Cette méthode est utilisée pour définir les données d'un étudiant. Elle prend en paramètre l'identifiant, le nom, le sexe, le total des notes et le pourcentage de l'élève et les affecte aux attributs correspondants de l'instance.
        showData(self) : Cette méthode est utilisée pour afficher les données de l'étudiant. Elle affiche l'ID, le nom, le sexe, le total des notes et le pourcentage de l'étudiant.
    Une instance de la classe Student est créée et nommée s.
    La méthode setData est appelée sur l'instance s pour définir les données de l'étudiant. Dans ce cas, l'ID de l'étudiant est fixé à 1, son nom à "Sam Kumar", son sexe à "Male", le total de ses notes à 422 et son pourcentage à 84,44.
    La méthode showData est appelée sur l'instance s pour imprimer les données de l'étudiant.
        \end{solution}
        

        \question
        Écrire un programme python pour créer une classe d'étudiants avec un constructeur et un destructeur.
        \par
        \begin{solution}
            \renewcommand{\nomfichier}{q711.py}
            \pythonfile{\chemincode \nomfichier}[][\nomfichier]
            Ce code Python définit une classe Student avec un constructeur, un destructeur et des méthodes pour définir et afficher les données des étudiants. Voici une explication du code :

    class Student: : Cette ligne définit une classe nommée Student, qui sert de modèle pour la création d'objets étudiants.
    def __init__(self): : Il s'agit de la méthode de construction de la classe Étudiant. Le constructeur est automatiquement appelé lorsqu'une instance de la classe est créée. À l'intérieur du constructeur :
        Des valeurs par défaut sont attribuées aux variables d'instance (attributs) telles que id, name, gender, total et per.
        Ces attributs représentent l'identifiant, le nom, le sexe, le total des notes et le pourcentage de l'étudiant.
    def __del__(self): : Il s'agit de la méthode de destruction de la classe Student. Le destructeur est appelé lorsqu'un objet est sur le point d'être détruit. Dans ce cas, il affiche simplement "Object Destroyed" lorsque l'objet est détruit.
    def setData(self): : Cette méthode permet de définir les données d'une instance d'étudiant. Elle demande à l'utilisateur de saisir des données et définit les valeurs des attributs de l'élève en fonction des données saisies par l'utilisateur.
    def showData(self): : Cette méthode affiche les données de l'élève. Elle imprime les valeurs des attributs de l'étudiant dans la console.
    s = Student() : Cette ligne crée une instance de la classe Student et l'affecte à la variable s. Cette instance représente un étudiant spécifique.
    s.setData() : La méthode setData est appelée sur l'instance s pour définir les données de cet étudiant. L'utilisateur est invité à saisir l'identifiant, le nom, le sexe, le total des notes et le pourcentage de l'élève.
    s.showData() : La méthode showData est appelée sur l'instance s pour afficher les données de l'élève. Elle imprime les valeurs des attributs définis dans la méthode setData.

Lorsque vous exécutez ce programme, il crée un objet Étudiant, définit ses attributs en fonction des données saisies par l'utilisateur, puis affiche les informations relatives à l'étudiant.
        \end{solution}
        

        \question
        Ecrire un programme python pour implémenter des Getters et des Setters dans une classe.
        \par
        \begin{solution}
            \renewcommand{\nomfichier}{q712.py}
            \pythonfile{\chemincode \nomfichier}[][\nomfichier]
            Ce programme Python définit une classe d'étudiant qui encapsule les détails de l'étudiant tels que le nom, le total et le pourcentage. Il démontre l'utilisation des méthodes setter et getter pour définir et récupérer ces attributs. Voici une explication du programme :

    classe Étudiant: : Il s'agit de la définition de la classe Student.
    def __init__(self): : Il s'agit de la méthode de construction de la classe Étudiant. Elle initialise trois variables d'instance : name, total et per avec des valeurs par défaut.
    self.name = "", self.total = "", self.per = 0 : Ces lignes initialisent le nom comme une chaîne vide, le total comme une chaîne vide et per comme 0.
    def setName(self, name): : Il s'agit d'une méthode setter permettant de définir le nom de l'élève.
    self.name = name : L'attribut name est mis à jour avec la valeur passée à la méthode setName.
    def getName(self): : Il s'agit d'une méthode getter permettant de récupérer le nom de l'élève.
    return self.name : L'attribut name est renvoyé lorsque la méthode getName est appelée.
    Des méthodes setter et getter similaires sont définies pour les attributs total et par.
    name = input("Enter a Name :") : Le programme invite l'utilisateur à saisir un nom et la saisie est stockée dans la variable name.
    total = int(input("Enter a Total :")) : Le programme invite l'utilisateur à saisir un total, et la saisie est convertie en un nombre entier et stockée dans la variable total.
    per = float(input("Enter a Percentage :")) : Le programme invite l'utilisateur à saisir un pourcentage, et la donnée est convertie en un nombre flottant et stockée dans la variable per.
    s = Student() : Une instance de la classe Student est créée, représentant un étudiant.
    s.setName(name), s.setTotal(total), s.setPercentage(per) : Les méthodes de définition sont utilisées pour définir le nom, le total et le pourcentage de l'étudiant en fonction des données fournies par l'utilisateur.
    n = s.getName(), t = s.getTotal(), p = s.getPercentage() : Les méthodes Getter sont utilisées pour récupérer le nom, le total et le pourcentage de l'élève.
    print("\nAffichage des détails de l'étudiant") : Un message est imprimé pour indiquer que les détails de l'élève vont être affichés.
    print("Nom :", n), print("Total :", t), print("Pourcentage :", p) : Le nom, le total et le pourcentage de l'étudiant sont affichés sur la console en utilisant les valeurs récupérées par les méthodes getter.

Ce programme démontre l'utilisation des principes orientés objet, de l'encapsulation et des méthodes setter et getter pour gérer et récupérer les informations relatives aux étudiants. Il permet à l'utilisateur de saisir des informations sur l'étudiant et affiche ensuite les détails de l'étudiant.
        \end{solution}
        

        \question
        Écrire un programme python pour mettre en œuvre l'abstraction à l'aide d'une classe abstraite.
        \par
        \begin{solution}
            \renewcommand{\nomfichier}{q713.py}
            \pythonfile{\chemincode \nomfichier}[][\nomfichier]
            Ce programme Python démontre l'utilisation de classes abstraites et de la superposition de méthodes dans la programmation orientée objet. Il définit une classe de base abstraite, Vehicle, et deux sous-classes concrètes, Bike et Car. La classe Vehicle contient deux méthodes abstraites, accelerate et park, qui sont surchargées par les classes enfants afin de fournir des implémentations spécifiques pour différents types de véhicules. Voici une explication du programme :

    classe Véhicule: : Il s'agit de la classe de base abstraite Véhicule. Elle contient deux méthodes abstraites, accelerate et park, qui sont définies mais n'ont pas d'implémentation (indiquée par l'instruction pass). Les méthodes abstraites sont des méthodes qui sont déclarées dans la classe de base mais qui n'ont pas d'implémentation spécifique, et elles doivent être surchargées par les classes enfants.
    classe Vélo(Véhicule): : Il s'agit de la classe Bike, qui hérite de la classe Vehicle. Elle surcharge les méthodes accelerate et park pour fournir un comportement spécifique à un vélo.
    classe Car(Véhicule): : Voici la classe Voiture, qui hérite également de la classe Véhicule. Elle surcharge les méthodes d'accélération et de stationnement pour fournir un comportement spécifique à une voiture.
    def accelerate(self, name): : Dans les classes Bike et Car, la méthode accelerate est implémentée pour accepter un paramètre name et imprimer un message spécifique au type de véhicule, indiquant qu'il accélère.
    def park(self, name): : Dans les classes Bike et Car, la méthode park est implémentée pour accepter un paramètre name et imprimer un message spécifique au type de véhicule, indiquant qu'il est en train de se garer.
    c = Voiture() : Une instance de la classe Voiture est créée et affectée à la variable c.
    c.accelerate("Voiture") : La méthode accelerate est appelée sur l'objet c, en passant "Car" comme paramètre de nom. Elle imprime un message indiquant que la voiture accélère.
    c.park("Voiture") : La méthode park est appelée sur l'objet c, en passant "Voiture" comme paramètre de nom. Elle imprime un message indiquant que la voiture est en train de se garer.
    b = Bike() : Une instance de la classe Bike est créée et assignée à la variable b.
    b.accelerate("Vélo") : La méthode accelerate est appelée sur l'objet b, en passant "Bike" comme paramètre de nom. Elle imprime un message indiquant que le vélo accélère.
    b.park("Vélo") : La méthode park est appelée sur l'objet b, en passant "Bike" comme paramètre de nom. Elle imprime un message indiquant que le vélo est en train de se garer.

Dans ce programme, l'utilisation d'une classe de base abstraite (Véhicule) garantit que toutes les sous-classes (Vélo et Voiture) fournissent des implémentations spécifiques pour les méthodes accelerate et park. Cela vous permet de créer différentes instances de véhicules (vélo et voiture) et d'appeler leurs méthodes spécifiques pour démontrer le comportement de chaque type de véhicule.
        \end{solution}
        

        \question
        Écrire un programme python pour implémenter une interface à l'aide d'une classe
        \par
        \begin{solution}
            \renewcommand{\nomfichier}{q714.py}
            \pythonfile{\chemincode \nomfichier}[][\nomfichier]
            Ce programme Python définit deux classes, Circle et Rectangle, qui héritent d'une classe de base commune, Shape. Le programme calcule les surfaces des cercles et des rectangles en fonction des données saisies par l'utilisateur. Voici une explication du programme :

    Importation du module mathématique : Le programme commence par importer le module mathématique pour accéder à des fonctions mathématiques telles que pi et pow.
    Classe Shape : Il s'agit de la classe de base, qui peut être considérée comme une interface dans ce contexte. Elle définit trois méthodes : input, process et output. Ces méthodes sont des espaces réservés pour des implémentations spécifiques dans les classes dérivées, mais ne sont pas implémentées ici.
    classe Circle(Shape) : Il s'agit de la classe Circle, qui hérite de la classe Shape. Elle représente un cercle et est responsable du calcul de sa surface.
        __init__(self, rad=0.0) : Le constructeur initialise l'objet Circle avec un rayon par défaut de 0,0. Il initialise également une variable de surface à 0,0.
        setdata(self) : Cette méthode invite l'utilisateur à saisir le rayon du cercle et le fixe dans l'attribut radius.
        circle_area(self) : Cette méthode calcule la surface du cercle à l'aide de la formule π * radius^2 et stocke le résultat dans l'attribut area.
        getdata(self) : Cette méthode imprime la surface calculée sur la console.
    classe Rectangle(Forme) : Il s'agit de la classe Rectangle, qui hérite de la classe Shape. Elle représente un rectangle et calcule sa surface.
        __init__(self, len=0, bre=0) : Le constructeur initialise l'objet Rectangle avec des valeurs de longueur et de largeur par défaut de 0. Il initialise également un attribut de surface à 0.
        setdata(self) : Cette méthode invite l'utilisateur à saisir la longueur et la largeur du rectangle et les définit dans les attributs respectifs.
        rect_area(self) : Cette méthode calcule la surface du rectangle en multipliant la longueur et la largeur et stocke le résultat dans l'attribut area.
        getdata(self) : Cette méthode imprime la surface calculée sur la console.
    Création d'instances de cercle et de rectangle :
        c = Circle() : Crée une instance de la classe Circle.
        c.setdata() : Invite l'utilisateur à saisir le rayon et le fixe.
        c.circle_area() : Calcule la surface du cercle.
        c.getdata() : Affiche la surface du cercle.
        r = Rectangle() : Crée une instance de la classe Rectangle.
        r.setdata() : Invite l'utilisateur à saisir la longueur et la largeur et les fixe.
        r.rect_area() : Calcule la surface du rectangle.
        r.getdata() : Affiche la surface du rectangle.

En résumé, ce programme montre comment les classes peuvent hériter d'une classe de base commune (ou d'une interface) et comment chaque classe dérivée fournit ses propres implémentations pour des méthodes spécifiques. Il calcule et affiche les surfaces des cercles et des rectangles en fonction des entrées de l'utilisateur.
        \end{solution}
        

        \question
        Écrire un programme pour l'héritage simple en Python
        \par
        \begin{solution}
            \renewcommand{\nomfichier}{q715.py}
            \pythonfile{\chemincode \nomfichier}[][\nomfichier]
            Ce programme Python démontre l'héritage dans la programmation orientée objet. Il définit deux classes, Détails et Étudiant, où la classe Étudiant hérite de la classe Détails. Voici une explication du programme :

    classe Détails: : Il s'agit de la classe Détails, qui sert de classe mère.
        __init__(self) : La méthode du constructeur initialise trois variables d'instance : name, total et per. Toutes ces variables sont initialement définies à 0 ou à une chaîne vide.
        setData(self, name, total, per) : Cette méthode permet de définir les valeurs des variables d'instance name, total et per.
        showData(self) : Cette méthode affiche les valeurs de name, total et per.
    class Student(Details): : La classe Student est une classe enfant qui hérite de la classe Details. Cela signifie qu'elle hérite des attributs et des méthodes de la classe mère.
        __init__(self) : La méthode du constructeur de la classe Student initialise une variable d'instance supplémentaire, Grade, qui prend initialement la forme d'une chaîne vide.
        setStudent(self, name, total, per, grade) : Cette méthode permet de définir les valeurs de name, total, per et de la nouvelle variable d'instance grade. Elle appelle la méthode setData de la classe mère pour définir les attributs communs.
        showStudent(self) : Cette méthode appelle d'abord la méthode showData de la classe mère pour afficher les détails hérités de Details. Elle affiche ensuite l'attribut grade.
    e = Student() : Une instance de la classe Student est créée et assignée à la variable e.
    e.setStudent("Kim", 430, 86.00, 'B') : La méthode setStudent est appelée sur l'objet e pour définir les détails de l'étudiant, y compris le nom, le total, le pourcentage et la note.
    e.showStudent() : La méthode showStudent est appelée pour afficher les détails de l'étudiant, notamment son nom, son total, son pourcentage et sa note.

Lorsque vous exécutez le programme, il crée une instance de la classe Student, définit les détails de l'étudiant à l'aide de la méthode setStudent, puis affiche ces détails à l'aide de la méthode showStudent. Ce programme illustre le concept d'héritage, dans lequel la classe Student hérite des attributs et des méthodes de la classe Details et les étend avec ses propres attributs et méthodes.
        \end{solution}
        

        \question
        Écrire un programme d'héritage avec deux classes enfantines (dérivées) en Python
        \par
        \begin{solution}
            \renewcommand{\nomfichier}{q716.py}
            \pythonfile{\chemincode \nomfichier}[][\nomfichier]
            Ce programme Python définit trois classes : Détails, Étudiant et Personnel. Ces classes sont utilisées pour collecter et afficher des informations sur les étudiants et les membres du personnel. Voici une explication du programme :

    classe Détails: : Il s'agit de la classe de base, qui contient des attributs et des méthodes communs dont héritent les classes Student et Staff.
        __init__(self) : Le constructeur initialise les attributs idn, name et gender aux valeurs par défaut.
        setDetails(self) : Cette méthode permet de définir les détails d'une personne, notamment son numéro d'identification, son nom et son sexe. Elle recueille ces informations auprès de l'utilisateur par le biais d'une entrée.
        showDetails(self) : Cette méthode permet d'afficher les détails d'une personne, notamment son numéro d'identification, son nom et son sexe.
    class Student(Details): : Il s'agit d'une sous-classe de la classe Details, spécialement conçue pour contenir des informations sur les étudiants.
        __init__(self) : Le constructeur initialise les attributs total et par aux valeurs par défaut. Il appelle également le constructeur de la classe mère (Details) en utilisant super().
        setStudent(self) : Cette méthode recueille des informations supplémentaires relatives aux étudiants, telles que le total des notes et le pourcentage, en utilisant les données fournies par l'utilisateur. Elle appelle également la méthode setDetails de la classe mère pour collecter les détails communs.
        showStudent(self) : Cette méthode affiche tous les détails d'un étudiant, y compris les détails communs (ID, nom et sexe) de la classe Details et les détails spécifiques (notes totales et pourcentage) de la classe Student.
    classe Staff(Details): : Il s'agit d'une autre sous-classe de la classe Détails, conçue pour contenir des informations sur les membres du personnel.
        __init__(self) : Le constructeur initialise les attributs depart et salary aux valeurs par défaut et appelle le constructeur de la classe mère (Details) à l'aide de super().
        setStaff(self) : Cette méthode recueille des informations supplémentaires relatives au personnel, telles que le département et le salaire, en utilisant les données fournies par l'utilisateur. Elle appelle également la méthode setDetails de la classe mère pour collecter les informations communes.
        showStaff(self) : Cette méthode affiche tous les détails d'un membre du personnel, y compris les détails communs (ID, nom et sexe) de la classe Details et les détails spécifiques (département et salaire) de la classe Staff.
    Le programme commence par créer des instances des classes Étudiant et Personnel afin de collecter et d'afficher les informations relatives à un étudiant et à un membre du personnel. Il demande à l'utilisateur de saisir les détails de chacun d'eux, puis affiche les informations.

Ce programme démontre l'utilisation de l'héritage de classe, où les classes Étudiant et Personnel héritent d'attributs et de méthodes communs de la classe Détails. Il montre comment vous pouvez collecter et afficher des informations spécifiques pour différents types d'objets (étudiants et membres du personnel) tout en réutilisant et en personnalisant les fonctionnalités communes fournies par la classe mère.
        \end{solution}
        

        \question
        Écrire un programme pour l'héritage multiple en Python
        \par
        \begin{solution}
            \renewcommand{\nomfichier}{q717.py}
            \pythonfile{\chemincode \nomfichier}[][\nomfichier]
            Ce programme Python démontre le concept d'héritage multiple en créant une classe enfant, Student, qui hérite de deux classes parents, PersonalInfo et AcademicInfo. Le programme recueille les informations personnelles et académiques de l'utilisateur et affiche ensuite les détails combinés de l'étudiant, y compris les informations personnelles et académiques. Voici une explication du programme :

    classe PersonalInfo : Il s'agit de la classe mère des informations personnelles. Elle contient un constructeur __init__ qui initialise les attributs tels que l'identifiant, le nom, le sexe, l'adresse et le contact. Elle comprend également une méthode display_personal_info qui permet d'afficher les informations personnelles.
    classe AcademicInfo : Il s'agit de la classe mère pour les informations académiques. Elle contient un constructeur __init__ qui initialise les attributs tels que la filière et l'année. Elle comprend également une méthode display_academic_info pour afficher les informations académiques.
    classe Student(PersonalInfo, AcademicInfo) : Il s'agit de la classe enfant qui hérite à la fois de PersonalInfo et de AcademicInfo. Elle surcharge le constructeur __init__ pour appeler les constructeurs des deux classes mères. Elle comprend également une méthode display_student_details pour afficher les informations personnelles et académiques.
    Dans la partie principale du programme, l'utilisateur est invité à saisir ses informations personnelles et académiques à l'aide de la fonction input.
    Une instance de la classe Student, nommée student, est créée en fournissant les informations collectées comme arguments au constructeur.
    Enfin, la méthode display_student_details est appelée sur l'objet étudiant pour afficher tous les détails de l'étudiant.

Lorsque vous exécutez ce programme, il collecte et affiche les informations personnelles et académiques d'un étudiant en héritant et en utilisant les méthodes et les attributs des classes mères PersonalInfo et AcademicInfo. Ceci démontre l'héritage multiple, où une classe enfant peut hériter de plusieurs classes parents pour réutiliser leurs attributs et méthodes.
        \end{solution}
        

        \question
        Écrire un programme python pour vérifier les nombres premiers en utilisant une approche orientée objet.
        \par
        \begin{solution}
            \renewcommand{\nomfichier}{q718.py}
            \pythonfile{\chemincode \nomfichier}[][\nomfichier]
            Ce programme Python définit une classe PrimeChecker qui peut être utilisée pour vérifier si un nombre donné est premier ou non. Voici une explication du fonctionnement du programme :

    La classe PrimeChecker est définie et prend un entier num comme paramètre lors de la création d'une instance de la classe.
    La méthode __init__ est le constructeur de la classe et initialise la variable d'instance self.num avec la valeur qui lui a été transmise.
    La méthode is_prime est utilisée pour déterminer si le nombre stocké dans self.num est premier ou non. Elle suit les étapes suivantes :
        Si le nombre est inférieur ou égal à 1, elle renvoie False, car les nombres premiers sont définis comme des entiers positifs supérieurs à 1.
        Si le nombre est égal à 2, il renvoie True, car 2 est le seul nombre premier pair.
        Si le nombre est pair (c'est-à-dire divisible par 2), il renvoie False car les nombres premiers (autres que 2) sont toujours impairs.
        Il vérifie ensuite la divisibilité du nombre à partir de 3 jusqu'à la racine carrée du nombre. Il effectue cette opération dans une boucle avec un pas de 2 pour ne vérifier que les nombres impairs, car les nombres pairs supérieurs à 2 ne peuvent pas être premiers. S'il trouve un diviseur dans cet intervalle, il renvoie False, indiquant que le nombre n'est pas premier.
        Si aucune des conditions ci-dessus n'est remplie, il renvoie True, ce qui indique que le nombre est premier.
    Le programme prend ensuite en compte la saisie d'un nombre par l'utilisateur à l'aide de input() et le convertit en un nombre entier à l'aide de int(). Ce nombre est stocké dans la variable num.
    Une instance de la classe PrimeChecker est créée avec les données de l'utilisateur et est affectée à la variable checker.
    La méthode is_prime de l'instance de checker est appelée. Si la méthode renvoie True, cela signifie que le nombre est premier et un message indiquant que le nombre est premier est imprimé. Si la méthode renvoie False, cela signifie que le nombre n'est pas premier et un message correspondant est imprimé.
        \end{solution}
        

        \question
        Écrire un programme python pour compter le nombre d'objets créés
        \par
        \begin{solution}
            \renewcommand{\nomfichier}{q719.py}
            \pythonfile{\chemincode \nomfichier}[][\nomfichier]
            Ce programme Python définit une classe appelée Student et démontre l'utilisation des variables de classe, de la création d'objets et des méthodes d'instance. Voici une explication du programme :

    classe Étudiant: : Cette ligne définit une classe nommée Student.
    count = 0 : Il s'agit d'une variable de classe nommée count qui est utilisée pour garder une trace du nombre d'objets Student créés. Elle est initialisée à 0.
    def __init__(self, name, age): : Il s'agit de la méthode de construction de la classe Étudiant. Elle est utilisée pour initialiser les attributs d'un objet Étudiant. self est une référence à l'instance de la classe en cours de création, et name et age sont les attributs qui sont définis pour chaque objet.
    self.name = name et self.age = age : Ces lignes définissent les attributs nom et âge de l'objet Étudiant avec les valeurs passées en argument.
    Student.count += 1 : Cette ligne incrémente la variable de classe count de 1 à chaque fois qu'un nouvel objet Etudiant est créé. Cela permet de garder une trace du nombre d'objets Étudiant créés.
    def GetDetails(self): : Il s'agit d'une méthode d'instance de la classe Étudiant, utilisée pour imprimer le nom et l'âge d'un étudiant.
    print("Name :", self.name) et print("Age :", self.age) : Ces lignes sont utilisées pour imprimer le nom et l'âge d'un objet Étudiant lorsque la méthode GetDetails est appelée.
    Quatre instances de la classe Étudiant sont créées :
        s1 = Student("Sam Kumar", 21) : un objet Student est créé avec le nom "Sam Kumar" et l'âge 21.
        s2 = Student("Tiya", 20) : un autre objet Student est créé avec le nom "Tiya" et l'âge de 20 ans.
        s3 = Student("Sathish", 19) : Un autre objet Student est créé avec le nom "Sathish" et l'âge 19.
        s3 = Étudiant("Deepika", 21) : Un autre objet Etudiant est créé avec le nom "Deepika" et l'âge de 21 ans. Notez que cette ligne réaffecte la variable s3, écrasant l'objet précédent créé avec cette variable.
    print("Nombre d'objets : ", Student.count) : Cette ligne imprime le nombre total d'objets Étudiant créés en accédant à la variable de classe count. Elle reflète le nombre d'objets créés, qui dans ce cas est de 3 puisque la variable s3 a été réaffectée.
        \end{solution}
        

        \question
        Ecrire un programme python pour vérifier le numéro Armstrong en utilisant une approche orientée objet.
        \par
        \begin{solution}
            \renewcommand{\nomfichier}{q720.py}
            \pythonfile{\chemincode \nomfichier}[][\nomfichier]
            Ce programme Python définit une classe appelée ArmstrongChecker qui vérifie si un nombre donné est un nombre Armstrong. Voici une explication pas à pas de son fonctionnement :

    La classe ArmstrongChecker est définie par un constructeur __init__, qui prend un argument, num, représentant le nombre à vérifier.
    Dans le constructeur, self.num se voit attribuer la valeur de l'argument num, qui stocke le nombre à vérifier pour être un nombre Armstrong.
    La classe possède une méthode appelée is_armstrong(). Cette méthode vérifie si le nombre stocké dans self.num est un nombre d'Armstrong et renvoie une valeur booléenne (True ou False).
    A l'intérieur de la méthode is_armstrong : a. num_str est calculé en convertissant le nombre en une chaîne de caractères. Cela permet de compter le nombre de chiffres dans le nombre original.
        num_digits est calculé comme la longueur de num_str, ce qui donne le nombre de chiffres dans le nombre original.
        digit_sum est calculé à l'aide d'une expression de générateur. Il calcule la somme de chaque chiffre élevé à la puissance de num_chiffres. Par exemple, si num est 153, digit_sum sera calculé comme suit : (1^3 + 5^3 + 3^3), soit 153.
        Enfin, la méthode renvoie True si digit_sum est égal au nombre original self.num, ce qui indique que le nombre est un nombre Armstrong. Dans le cas contraire, elle renvoie False.
    Le programme prend ensuite en compte la saisie d'un nombre par l'utilisateur en utilisant int(input("Enter a Number : ")) et le stocke dans la variable num.
    Une instance de la classe ArmstrongChecker est créée avec le nombre fourni num comme argument, et l'instance est stockée dans la variable checker.
    Le programme vérifie ensuite si le nombre est un nombre Armstrong en appelant la méthode is_armstrong() de l'objet checker.
    Si le nombre est un nombre d'Armstrong (c'est-à-dire que la méthode renvoie True), il imprime un message indiquant que le nombre est un nombre d'Armstrong. Dans le cas contraire, un message indiquant que le nombre n'est pas un nombre d'Armstrong s'affiche.

Le programme vérifie si un nombre donné est un nombre d'Armstrong en utilisant la définition du nombre d'Armstrong (un nombre est un nombre d'Armstrong si la somme de ses chiffres élevée à la puissance du nombre de chiffres est égale au nombre original). Pour ce faire, il encapsule la logique dans une classe appelée ArmstrongChecker.
        \end{solution}
        

        \question
        Écrire un programme python pour l'héritage multi-niveaux
        \par
        \begin{solution}
            \renewcommand{\nomfichier}{q721.py}
            \pythonfile{\chemincode \nomfichier}[][\nomfichier]
            Ce programme Python définit une hiérarchie de classes avec un héritage à plusieurs niveaux en Python. Voici un aperçu des classes et de leurs relations :

    Personne (classe de base) : Cette classe possède un constructeur __init__ qui initialise les attributs nom et âge. Elle possède également une méthode display_info qui imprime le nom et l'âge.
    Étudiant (classe intermédiaire) : Cette classe hérite de Person et ajoute un nouvel attribut student_id. Elle possède son propre constructeur, qui prend en paramètre le nom, l'âge et l'identifiant de l'étudiant. Le constructeur de la classe Student appelle le constructeur de la classe Person en utilisant super() pour définir le nom et l'âge. Elle possède également une méthode display_student_info, qui appelle la méthode display_info de la classe Person et ajoute les informations relatives à l'identifiant de l'étudiant.
    GraduateStudent (classe dérivée) : Cette classe hérite de la classe Student et ajoute un nouvel attribut research_topic. Elle possède son propre constructeur, qui prend en paramètre le nom, l'âge, l'identifiant de l'étudiant et le sujet de la recherche. Le constructeur de la classe GraduateStudent appelle le constructeur de la classe Student en utilisant super() pour définir le nom, l'âge et l'identifiant de l'étudiant. Il possède également une méthode display_graduate_info, qui appelle la méthode display_student_info de la classe Student et ajoute les informations research_topic.

Vous créez une instance de la classe GraduateStudent avec le nom "Alice", l'âge de 25 ans, le numéro d'étudiant "STU001" et le sujet de recherche "Machine Learning". Vous appelez ensuite la méthode display_graduate_info sur l'objet graduate_student, qui imprime toutes les informations de la classe de base Person, de la classe Student et de la classe GraduateStudent, y compris le sujet de recherche.
        \end{solution}
        

        \question
        Ecrire un programme python pour vérifier le nombre de palindromes en utilisant une approche orientée objet.
        \par
        \begin{solution}
            \renewcommand{\nomfichier}{q722.py}
            \pythonfile{\chemincode \nomfichier}[][\nomfichier]
            Le programme Python PalindromeChecker est conçu pour vérifier si un nombre donné est un palindrome. Voici une explication du fonctionnement du programme :

    class PalindromeChecker : Il s'agit d'une classe qui définit un vérificateur de palindromes. Elle possède deux méthodes :
        __init__(self, num) : Il s'agit de la méthode du constructeur qui initialise l'objet avec le nombre num.
        is_palindrome(self) : Cette méthode vérifie si le nombre est un palindrome. Pour ce faire, elle convertit d'abord le nombre en une chaîne de caractères (num_str). Ensuite, elle inverse la chaîne en utilisant le découpage (reversed_str = num_str[::-1] ) et vérifie si la chaîne inversée est égale à la chaîne originale. Si c'est le cas, il renvoie True, ce qui indique que le nombre est un palindrome ; sinon, il renvoie False.
    num = int(input("Enter a Number : ")) : Cette ligne de code invite l'utilisateur à saisir un nombre et convertit la saisie de l'utilisateur en un nombre entier, qui est stocké dans la variable num.
    checker = PalindromeChecker(num) : Cette ligne crée une instance de la classe PalindromeChecker avec le nombre saisi par l'utilisateur. La valeur num est transmise comme argument au constructeur, initialisant l'attribut num de l'objet checker.
    if checker.is_palindrome(): : Cette instruction conditionnelle vérifie si la méthode is_palindrome de l'objet checker renvoie True. Si c'est le cas, cela signifie que le nombre saisi est un palindrome et un message indiquant que le nombre est un palindrome s'affiche. Dans le cas contraire, un message indiquant que le nombre n'est pas un palindrome est imprimé.
    Le programme imprime alors "{num} est un nombre palindrome" ou "{num} n'est pas un nombre palindrome" en fonction du résultat de la méthode is_palindrome.

Voici comment fonctionne le programme :

    L'utilisateur est invité à saisir un nombre.
    Le programme crée une instance de la classe PalindromeChecker avec le nombre saisi.
    Il vérifie si le nombre saisi est un palindrome à l'aide de la méthode is_palindrome.
    En fonction du résultat, il imprime un message indiquant si le nombre est un palindrome ou non.

Ce programme démontre le concept d'utilisation d'une classe pour encapsuler la fonctionnalité permettant de vérifier si un nombre est un palindrome, ce qui rend le code plus organisé et réutilisable.
        \end{solution}
        

        \question
        Écrire un programme pour le programme d'enregistrement de la taille des étudiants pour une école en Python
        \par
        \begin{solution}
            \renewcommand{\nomfichier}{q723.py}
            \pythonfile{\chemincode \nomfichier}[][\nomfichier]
            Le code Python définit une classe appelée SchoolHeightRecords qui vous permet de gérer les enregistrements de taille des élèves. Voici la décomposition du code :

    La classe SchoolHeightRecords est définie avec les méthodes suivantes :
        __init__ : Initialise un dictionnaire vide appelé student_records pour stocker les noms des élèves et leurs tailles correspondantes.
        add_student : Ajoute la taille d'un nouvel élève au dictionnaire student_records.
        remove_student : Supprime l'enregistrement d'un élève du dictionnaire student_records en fournissant le nom de l'élève.
        find_student_height : Permet d'obtenir la taille d'un élève en fournissant son nom. Si l'élève n'est pas trouvé, cette fonction renvoie None.
        display_records : Affiche tous les enregistrements de la taille des élèves stockés dans le dictionnaire.
    Une instance de la classe SchoolHeightRecords est créée avec la variable school_records.
    Un menu est affiché à l'utilisateur avec les options suivantes :
        1. Ajouter la taille de l'élève
        2. Supprimer la taille de l'élève
        3. Trouver la taille d'un élève
        4. Afficher tous les enregistrements
        5. Quitter
    Le programme entre dans une boucle où l'utilisateur est invité à saisir son choix.
    En fonction du choix de l'utilisateur, la méthode correspondante de l'instance SchoolHeightRecords est appelée pour effectuer l'opération souhaitée. L'utilisateur peut ajouter, supprimer, rechercher ou afficher des relevés de taille d'élèves, ou quitter le programme.
    Si l'utilisateur fournit un choix non valide, il en est informé par le message "Choix non valide" et est invité à réessayer.

Le code vous permet de gérer de manière interactive les enregistrements de taille des élèves et il continuera à fonctionner jusqu'à ce que l'utilisateur choisisse de quitter le programme (option 5).

Il convient de noter que les tailles sont stockées sous forme de chaînes de caractères dans le dictionnaire student_records. Si vous devez effectuer des opérations numériques avec les tailles, vous devez les convertir en entiers ou en flottants avant de les utiliser dans les calculs.
        \end{solution}
        

        \question
        Ecrire un programme python pour gérer l'enregistrement d'un magasin de téléphonie (mobile shop) en utilisant la classe
        \par
        \begin{solution}
            \renewcommand{\nomfichier}{q724.py}
            \pythonfile{\chemincode \nomfichier}[][\nomfichier]
            Le code Python définit deux classes : Phone et PhoneStore, qui permettent de gérer l'inventaire des téléphones d'un magasin. Voici la décomposition du code :

    La classe Phone représente des téléphones individuels. Chaque téléphone possède des attributs pour sa marque, son modèle et son prix. La méthode __init__ initialise ces attributs lorsqu'un nouvel objet téléphone est créé.
    La classe PhoneStore représente l'inventaire des téléphones d'un magasin. Elle possède les méthodes suivantes :
        __init__ : Initialise une liste vide appelée inventaire pour stocker les objets téléphone.
        add_phone : Ajoute un objet téléphone à la liste d'inventaire.
        remove_phone : Supprime un téléphone de l'inventaire en fonction de sa marque et de son modèle. Si le téléphone est trouvé et supprimé, un message de réussite est affiché ; sinon, un message "not found" (introuvable) est affiché.
        find_phone : Recherche un téléphone dans l'inventaire en fonction de sa marque et de son modèle. S'il est trouvé, il renvoie l'objet téléphone ; sinon, il renvoie None.
        display_inventory : Affiche l'inventaire actuel des téléphones, y compris la marque, le modèle et le prix.
    Une instance de la classe PhoneStore est créée avec la variable phone_store.
    Un menu est affiché à l'utilisateur avec les options suivantes :
        1. Ajouter un téléphone à l'inventaire
        2. Supprimer un téléphone de l'inventaire
        3. Trouver un téléphone dans l'inventaire
        4. Afficher l'inventaire
        5. Quitter
    Le programme entre dans une boucle où l'utilisateur est invité à saisir son choix (sous forme de chaîne de caractères).
    En fonction du choix de l'utilisateur, la méthode correspondante de l'instance PhoneStore est appelée pour effectuer l'opération souhaitée. L'utilisateur peut ajouter, supprimer, rechercher ou afficher l'inventaire des téléphones, ou quitter le programme.
    Si l'utilisateur fournit un choix non valide, il est averti par le message "Choix non valide" et invité à réessayer.

Le code convertit les entrées de marque et de modèle en majuscules à l'aide de la méthode .capitalize() avant de les stocker dans les objets téléphone et lors de la recherche dans l'inventaire. Cela permet de s'assurer que la correspondance entre les majuscules et les minuscules est utilisée lors de l'ajout, de la suppression ou de la recherche de téléphones.

Le code vous permet de gérer de manière interactive l'inventaire des téléphones d'un magasin et continuera à fonctionner jusqu'à ce que l'utilisateur choisisse de quitter (option 5).
        \end{solution}
        

        \question
        Écrire un programme python pour additionner deux distances en utilisant les concepts de classe et d'objet
        \par
        \begin{solution}
            \renewcommand{\nomfichier}{q725.py}
            \pythonfile{\chemincode \nomfichier}[][\nomfichier]
            Ce programme Python définit une classe Distance qui représente une distance en kilomètres, mètres et centimètres. La classe fournit des méthodes permettant d'additionner deux distances et d'afficher une distance dans un format lisible par l'homme. Voici une explication du programme :

    classe Distance : Il s'agit de la classe principale représentant une distance. Elle possède les méthodes suivantes :
        __init__(self, km=0, m=0, cm=0) : Le constructeur initialise l'objet avec des valeurs optionnelles pour les kilomètres (0 par défaut), les mètres (0 par défaut) et les centimètres (0 par défaut).
        add(self, other_distance) : Cette méthode prend en paramètre un autre objet Distance, ajoute les composantes correspondantes (kilomètres, mètres et centimètres) et renvoie un nouvel objet Distance.
        display(self) : Cette méthode renvoie une chaîne de caractères formatée représentant la distance.
    Le programme invite ensuite l'utilisateur à saisir deux distances : une pour la distance1 et une pour la distance2. L'utilisateur est invité à saisir la distance en kilomètres, en mètres et en centimètres pour chacune d'entre elles.
    Les objets Distance distance1 et distance2 sont créés avec les valeurs fournies par l'utilisateur.
    La méthode add est appelée sur distance1, en passant distance2 comme paramètre. Il en résulte un nouvel objet Distance result_distance qui représente la somme des deux distances d'entrée.
    Le programme affiche le résultat en appelant la méthode display sur distance_résultat, qui fournit la somme des deux distances dans un format lisible par l'homme.

Lorsque vous exécutez ce programme, il vous permet de saisir deux distances, puis de calculer et d'afficher la somme de ces distances. Le programme se charge d'ajuster les unités si nécessaire, en veillant à ce que le résultat soit affiché correctement en termes de kilomètres, de mètres et de centimètres.
        \end{solution}
        

        \question
        Ecrire un programme python pour trouver la personne la plus âgée de deux personnes en utilisant la classe et l'objet.
        \par
        \begin{solution}
            \renewcommand{\nomfichier}{q726.py}
            \pythonfile{\chemincode \nomfichier}[][\nomfichier]
            Ce programme Python définit une classe Person et une fonction find_elder pour trouver la personne la plus âgée parmi deux individus donnés en fonction de leur âge. Voici une explication du programme :

    classe Personne : Il s'agit de la classe Personne avec un constructeur __init__ qui initialise deux attributs, le nom et l'âge. Elle représente une personne avec un nom et un âge.
    def find_elder(personne1, personne2) : Cette fonction prend deux objets Person, person1 et person2, comme paramètres. Elle compare leurs âges et renvoie la personne la plus âgée. Si les deux personnes ont le même âge, elle renvoie None.
    Le programme crée deux objets Personne, personne1 et personne2, représentant deux individus avec leurs noms et âges respectifs.
    La fonction find_elder est appelée avec ces deux objets Personne et le résultat est stocké dans la variable elder.
    Le programme vérifie ensuite si elder est None. Si c'est le cas, cela signifie que les deux personnes ont le même âge, et il imprime donc "Les deux personnes ont le même âge". Sinon, il imprime le nom et l'âge de la personne la plus âgée.

Lorsque vous exécutez ce programme, il détermine laquelle des deux personnes est la plus âgée en fonction de leur âge et imprime le message correspondant. Si les deux personnes ont le même âge, il indiquera que les deux ont le même âge. Ce programme montre comment définir une classe, créer des objets et comparer leurs attributs pour prendre des décisions basées sur l'état de l'objet.
        \end{solution}
        

        \question
        Écrire un programme python pour le système de gestion bancaire
        \par
        \begin{solution}
            \renewcommand{\nomfichier}{q727.py}
            \pythonfile{\chemincode \nomfichier}[][\nomfichier]
            Ce programme Python simule un système de gestion bancaire simple en utilisant la classe Bank. Il permet aux utilisateurs de créer des comptes, de déposer de l'argent, de retirer de l'argent, de vérifier les soldes et de quitter le programme. Voici une explication du programme :

    classe Banque : Il s'agit de la classe principale représentant la banque. Elle possède les méthodes suivantes :
        __init__(self) : Le constructeur initialise un dictionnaire de comptes vide pour stocker les informations relatives aux comptes.
        create_account(self, account_number, account_holder, initial_balance) : Cette méthode crée un nouveau compte et le stocke dans le dictionnaire des comptes. Elle vérifie qu'il n'y a pas de duplication de numéro de compte et que le solde initial n'est pas négatif.
        deposit(self, account_number, amount) : Cette méthode permet d'effectuer des dépôts sur un compte existant, en mettant à jour le solde du compte.
        withdraw(self, account_number, amount) : Cette méthode permet d'effectuer des retraits sur un compte existant, en mettant à jour le solde du compte. Elle vérifie également que le solde est suffisant et que les montants des retraits ne sont pas négatifs.
        check_balance(self, account_number) : Cette méthode vérifie le solde d'un compte existant et renvoie le nom et le solde du titulaire du compte sous la forme d'une chaîne formatée.
    Le programme crée un objet Banque nommé banque.
    Le programme présente un menu à l'utilisateur avec les options suivantes : créer un compte, déposer de l'argent, retirer de l'argent, vérifier le solde et quitter.
    Dans une boucle, le programme lit le choix de l'utilisateur et effectue l'opération sélectionnée :
        Pour "Créer un compte", il recueille les détails du compte et appelle la méthode create_account.
        Pour "Déposer de l'argent", il recueille le numéro de compte et le montant du dépôt et appelle la méthode de dépôt.
        Pour "Retirer de l'argent", il collecte le numéro de compte et le montant du retrait et appelle la méthode de retrait.
        Pour "Vérifier le solde", il collecte le numéro de compte et appelle la méthode check_balance.
        Pour "Quitter", il quitte le programme.
    Le programme assure la validation des entrées et fournit un retour d'information approprié à l'utilisateur.

Lorsque vous exécutez ce programme, il vous permet de simuler un système de gestion bancaire de base en créant des comptes, en déposant de l'argent, en retirant de l'argent et en vérifiant les soldes. Il utilise la classe Banque pour gérer les comptes et les opérations sur les comptes.
        \end{solution}
        

        \question
        Ecrire un programme python pour obtenir les détails de l'étudiant en entrée et imprimer le résultat après la mise à jour des notes.
        \par
        \begin{solution}
            \renewcommand{\nomfichier}{q728.py}
            \pythonfile{\chemincode \nomfichier}[][\nomfichier]
            Le programme Python pour un système simple de gestion des étudiants. Il vous permet de saisir les coordonnées des étudiants, de mettre à jour les notes des étudiants et d'imprimer les coordonnées des étudiants. Voici un aperçu du fonctionnement du code :

    classe Student : Cette classe définit une classe nommée Student, qui représente le système de gestion des étudiants.
    __init__(self) : Le constructeur initialise un dictionnaire vide appelé student_details en tant que variable d'instance.
    @property input_student_details(self) : Il s'agit d'une méthode de propriété qui vous permet de saisir les détails de l'étudiant. Elle invite l'utilisateur à saisir le numéro de rôle, le nom de l'étudiant et ses notes. Les détails sont stockés dans le dictionnaire student_details sous la clé roll number.
    update_student_marks(self) : Cette méthode permet de mettre à jour les notes d'un élève. Elle invite l'utilisateur à saisir le numéro de rôle de l'élève à mettre à jour, puis les nouvelles notes. Si l'élève est trouvé dans le dictionnaire student_details, ses notes sont mises à jour.
    print_student_details(self) : Cette méthode permet d'imprimer les détails d'un étudiant spécifique. Elle invite l'utilisateur à saisir le numéro de rôle de l'étudiant et, si l'étudiant est trouvé, son numéro de rôle, son nom et ses notes sont imprimés.
    Le programme crée une instance de la classe Étudiant appelée obj.
    Le code entre dans une boucle while qui affiche un menu avec quatre options : saisir les coordonnées de l'élève, mettre à jour les notes de l'élève, imprimer les coordonnées de l'élève et quitter.
    L'utilisateur est invité à saisir son choix (1 à 4) et, en fonction de ce choix, l'une des méthodes de l'instance obj est appelée.
    Si l'utilisateur choisit l'option 4, le programme sort de la boucle et se termine.
    Si l'utilisateur saisit un choix non valide, un message d'erreur s'affiche.
        \end{solution}
        

        \question
        Ecrire un programme python pour les tableaux d'objets
        \par
        \begin{solution}
            \renewcommand{\nomfichier}{q729.py}
            \pythonfile{\chemincode \nomfichier}[][\nomfichier]
            Le programme Python définit une classe d'étudiants et crée une liste (étudiants) pour stocker les instances de la classe d'étudiants. Il montre également comment accéder aux attributs d'un étudiant et les modifier. Voici la décomposition du code :

    classe Étudiant : Cette classe définit une classe appelée Student, qui représente un étudiant avec des attributs tels que le numéro de rôle, le nom et les notes.
    __init__(self, roll_number, name, marks) : La méthode du constructeur initialise un objet Student avec le numéro de rôle, le nom et les notes fournis.
    __str__(self) : Cette méthode spéciale renvoie une représentation sous forme de chaîne de caractères de l'objet Étudiant, affichant le numéro de rôle, le nom et les notes.
    étudiants = [] : Cette méthode crée une liste vide pour stocker les objets Étudiant.
    Trois étudiants sont créés et ajoutés à la liste des étudiants à l'aide de la méthode append. Chaque étudiant a un numéro de rôle, un nom et des notes uniques.
    Une boucle for itère sur la liste des étudiants et imprime les détails de chaque étudiant à l'aide de la méthode __str__.
    Elle montre ensuite comment accéder aux attributs d'un étudiant spécifique et les modifier. Dans ce cas, il sélectionne le premier étudiant de la liste, modifie son nom et ses notes, et imprime les informations mises à jour.
        \end{solution}
        

        \question
        Ecrire un programme python pour rechercher des objets dans un tableau d'objets en utilisant l'ID
        \par
        \begin{solution}
            \renewcommand{\nomfichier}{q730.py}
            \pythonfile{\chemincode \nomfichier}[][\nomfichier]
            Le programme définit une classe Étudiant, crée une liste d'objets Étudiant et vous permet de rechercher un étudiant par son numéro d'identification. Voici la décomposition du code :

    classe Étudiant : Cette classe représente un étudiant avec des attributs tels que l'identifiant de l'étudiant, le nom et le pourcentage.
    __init__(self, student_id, name, per) : Le constructeur initialise un objet Student avec l'identifiant, le nom et le pourcentage fournis.
    __str__(self) : Cette méthode spéciale renvoie une représentation sous forme de chaîne de caractères de l'objet Étudiant, affichant l'identifiant, le nom et le pourcentage de l'étudiant.
    Un tableau (liste) d'objets Étudiant est créé et initialisé avec les données de l'étudiant.
    search_student_by_id(student_id) : Cette fonction permet de rechercher un étudiant à partir de son numéro d'identification. Elle parcourt la liste des étudiants et renvoie le premier objet étudiant correspondant à l'identifiant spécifié, ou Aucun si aucune correspondance n'est trouvée.
    Le code imprime tous les identifiants des étudiants de la liste à l'aide d'une boucle for.
    L'utilisateur est invité à saisir l'identifiant de l'étudiant à rechercher.
    La fonction search_student_by_id est appelée avec l'identifiant fourni par l'utilisateur et, si un étudiant correspondant à cet identifiant est trouvé, il est imprimé avec un message "Étudiant trouvé". Si aucun étudiant correspondant n'est trouvé, un message "Étudiant non trouvé" s'affiche.
        \end{solution}
        

        \question
        Écrire un programme python pour l'initialisation du constructeur
        \par
        \begin{solution}
            \renewcommand{\nomfichier}{q731.py}
            \pythonfile{\chemincode \nomfichier}[][\nomfichier]
            Le programme définit une classe Étudiant et crée trois instances de la classe, chacune représentant un étudiant différent. Il affiche ensuite les détails de chaque étudiant à l'aide de la méthode display_details. Voici la décomposition du code :

    classe Étudiant : Cette classe définit un étudiant avec des attributs tels que le numéro de rôle, le nom et le pourcentage.
    __init__(self, roll_number, name, percent) : Le constructeur initialise un objet Student avec le numéro de rôle, le nom et le pourcentage fournis.
    display_details(self) : Cette méthode permet d'afficher les détails d'un étudiant, y compris son numéro de rôle, son nom et son pourcentage.
    Trois instances de la classe Étudiant (s1, s2 et s3) sont créées et initialisées avec des données différentes pour chaque étudiant.
    Le code affiche ensuite les détails de chaque étudiant en appelant la méthode display_details pour chaque instance.

Voici ce que fait le code :

    Il crée trois objets étudiants avec des données différentes.
    Il imprime les détails de chaque élève à l'aide de la méthode display_details, en affichant leur numéro de rôle, leur nom et leur pourcentage.
        \end{solution}
        

        \question
        Ecrire un programme python pour paramétrer un constructeur et un destructeur.
        \par
        \begin{solution}
            \renewcommand{\nomfichier}{q732.py}
            \pythonfile{\chemincode \nomfichier}[][\nomfichier]
            Le programme définit une classe Étudiant avec un constructeur et un destructeur. Il crée trois instances de la classe, affiche leurs détails, puis supprime explicitement ces instances pour démontrer la fonctionnalité du destructeur. Voici la décomposition du code :

    classe Étudiant : Cette classe définit un étudiant avec des attributs tels que le numéro de rôle, le nom et le pourcentage.
    __init__(self, roll_number, name, per) : Le constructeur initialise un objet Student avec le numéro de rôle, le nom et le pourcentage fournis.
    display_details(self) : Cette méthode est utilisée pour imprimer les détails d'un étudiant, y compris son numéro de rôle, son nom et son pourcentage.
    del__(self) : Cette méthode spéciale est le destructeur. Elle est automatiquement appelée lorsqu'une instance de la classe est supprimée. Dans ce cas, elle imprime un message indiquant le numéro de rôle de l'étudiant supprimé.
    Trois instances de la classe Étudiant (s1, s2 et s3) sont créées et initialisées avec des données différentes pour chaque étudiant.
    Le code affiche les détails de chaque étudiant en appelant la méthode display_details pour chaque instance.
    Les instances sont explicitement supprimées à l'aide de l'instruction del, qui déclenche la méthode du destructeur __del__.

Voici ce que fait le code :

    Il crée trois objets étudiants avec des détails différents.
    Il imprime les détails de chaque étudiant à l'aide de la méthode display_details.
    Il supprime explicitement chaque objet élève, ce qui déclenche le destructeur et imprime un message indiquant que l'élève est supprimé.
        \end{solution}
        

        \question
        Ecrire un programme python pour ajouter des objets '+' en utilisant l'opérateur
        \par
        \begin{solution}
            \renewcommand{\nomfichier}{q733.py}
            \pythonfile{\chemincode \nomfichier}[][\nomfichier]
            Le programme python définit une classe ComplexNumber qui représente les nombres complexes. Elle permet de créer des objets de type nombre complexe, de les afficher sous forme de chaînes de caractères et d'effectuer l'addition de deux nombres complexes à l'aide de l'opérateur +. Voici la décomposition du code :

    classe ComplexNumber : Cette classe représente un nombre complexe avec les attributs real et imag, qui représentent respectivement les parties réelle et imaginaire.
    __init__(self, real, imag) : Le constructeur initialise un objet ComplexNumber avec les parties réelle et imaginaire fournies.
    __str__(self) : Cette méthode spéciale renvoie une représentation sous forme de chaîne de caractères du nombre complexe au format "a + bi" où "a" est la partie réelle, "b" la partie imaginaire et "i" l'unité imaginaire.
    __add__(self, other) : Cette méthode spéciale surcharge l'opérateur +, ce qui permet d'additionner deux nombres complexes. Elle additionne les parties réelle et imaginaire séparément et renvoie un nouvel objet ComplexNumber représentant la somme.
    Deux objets de nombres complexes, complex1 et complex2, sont créés avec des parties réelles et imaginaires différentes.
    Le code utilise l'opérateur + pour additionner complex1 et complex2, ce qui produit un nouvel objet ComplexNumber appelé result.
    L'instruction print affiche le résultat de l'addition, qui est la somme des nombres complexes.

Voici ce que fait le code :

    Il crée deux objets nombres complexes avec des parties réelles et imaginaires différentes.
    Il additionne ces nombres complexes à l'aide de l'opérateur +, qui appelle la méthode __add__.
    Il imprime le résultat de l'addition.
        \end{solution}
        

        \question
        Ecrire un programme python pour comparer deux objets en utilisant l'opérateur '>'.
        \par
        \begin{solution}
            \renewcommand{\nomfichier}{q734.py}
            \pythonfile{\chemincode \nomfichier}[][\nomfichier]
            Le programme python définit une classe Number qui représente un nombre. Il permet de créer des objets Number, de les afficher sous forme de chaînes et d'effectuer une comparaison supérieure à l'aide de l'opérateur >. Voici la décomposition du code :

    classe Number : Cette classe représente un nombre avec une valeur d'attribut.
    __init__(self, value) : Le constructeur initialise un objet Number avec la valeur fournie.
    __str__(self) : Cette méthode spéciale renvoie une représentation sous forme de chaîne de caractères du nombre.
    __gt__(self, other) : Cette méthode spéciale surcharge l'opérateur > (plus grand que), ce qui permet de comparer deux objets Number. Elle vérifie si la valeur de l'objet courant est supérieure à la valeur de l'autre objet.
    Deux objets Number, number1 et number2, sont créés avec des valeurs différentes.
    Le code compare number1 et number2 en utilisant l'opérateur > dans une expression et stocke le résultat dans la variable res.
    En fonction du résultat, le code imprime un message indiquant si nombre1 est supérieur à nombre2.

Voici ce que fait le code :

    Il crée deux objets Number avec des valeurs différentes.
    Il compare les deux nombres à l'aide de l'opérateur >.
    Il imprime un message basé sur le résultat de la comparaison.
        \end{solution}
        

        \question
        Écrire un programme python pour passer des objets en tant qu'arguments et renvoyer des objets à partir d'une fonction
        \par
        \begin{solution}
            \renewcommand{\nomfichier}{q735.py}
            \pythonfile{\chemincode \nomfichier}[][\nomfichier]
            Le programme python définit une classe MyClass et une fonction process_object. La fonction prend en argument un objet de la classe MyClass, modifie l'objet en doublant sa valeur et renvoie l'objet modifié. Voici la décomposition du code :

    classe MaClasse : Cette classe représente un objet avec une valeur d'attribut.
    __init__(self, value) : Le constructeur initialise un objet MyClass avec la valeur fournie.
    def process_object(obj) : Cette fonction prend un objet obj comme argument, double l'attribut value de l'objet, puis renvoie l'objet modifié.
    Une instance de MyClass est créée et my_obj est initialisé avec une valeur de 10.
    La fonction process_object est appelée avec mon_obj comme argument. Elle modifie mon_obj en doublant son attribut value, et l'objet modifié est renvoyé et assigné à result_obj.
    Le code imprime l'attribut de valeur de l'objet original (mon_obj) et de l'objet modifié (résultat_obj).

Voici ce que fait le code :

    Il crée un objet MyClass avec une valeur initiale de 10.
    Il appelle la fonction process_object, qui modifie l'objet en doublant sa valeur.
    Il imprime l'attribut value de l'objet original et de l'objet modifié.
        \end{solution}
        

        \question
        Ecrire un programme python pour illustrer le fonctionnement des décorateurs
        \par
        \begin{solution}
            \renewcommand{\nomfichier}{q736.py}
            \pythonfile{\chemincode \nomfichier}[][\nomfichier]
            L'utilisation d'un décorateur en Python. Un décorateur est une fonction qui enveloppe une autre fonction pour lui ajouter des fonctionnalités. Dans ce code, vous avez défini une fonction décorateur my_decorator, qui enveloppe la fonction msg. Voici la décomposition du code :

    def my_decorator(func) : Il s'agit de la fonction décoratrice. Elle prend une fonction func en argument et définit une fonction interne appelée wrapper. La fonction wrapper ajoute un comportement avant et après l'appel de func.
    def wrapper() : Il s'agit de la fonction interne du décorateur. Elle ajoute un comportement avant et après l'appel de la fonction originale func.
    À l'intérieur de la fonction wrapper, des instructions d'impression indiquent que quelque chose se passe avant et après l'appel de la fonction d'origine.
    return wrapper : La fonction decorator renvoie la fonction wrapper.
    @my_decorator : Il s'agit d'une syntaxe de décorateur, indiquant que la fonction msg est décorée avec mon_décorateur. Cela signifie que lorsque vous appelez msg, elle sera enveloppée par la fonction wrapper définie dans mon_décorateur.
    def msg() : Il s'agit de la fonction à décorer. Elle imprime simplement "Hello world !".
    Lorsque vous appelez msg(), vous appelez en fait la fonction décorée. Le décorateur my_decorator enveloppe msg, de sorte qu'il imprime les messages supplémentaires avant et après l'impression de "Hello world !

La fonction msg est décorée par le décorateur my_decorator, de sorte que lorsque vous appelez msg(), elle affiche "Quelque chose se passe avant l'appel de la fonction", puis "Hello world !", et enfin "Quelque chose se passe après l'appel de la fonction". Cela montre comment les décorateurs peuvent ajouter un comportement aux fonctions sans modifier leur code source.
        \end{solution}
        

        \question
        Ecrire un programme python pour illustrer le fonctionnement de la méthode abstraite
        \par
        \begin{solution}
            \renewcommand{\nomfichier}{q737.py}
            \pythonfile{\chemincode \nomfichier}[][\nomfichier]
            L'utilisation des classes de base abstraites et de l'héritage en Python. Voici la décomposition du code :

    from abc import ABC, abstractmethod : Cette ligne importe la classe ABC (Abstract Base Class) et le décorateur abstractmethod du module abc. La classe ABC est utilisée pour créer des classes de base abstraites, et le décorateur abstractmethod est utilisé pour déclarer des méthodes abstraites dans ces classes.
    classe Shape(ABC) : Cette classe définit une classe de base abstraite Shape qui hérite de la classe ABC. Elle possède une zone de méthodes abstraites.
    @abstractmethod : Ce décorateur est utilisé pour déclarer la méthode area comme une méthode abstraite. Les méthodes abstraites doivent être implémentées par toute sous-classe concrète de la classe Shape.
    classe Circle(Forme) : Cette classe définit une sous-classe Circle de la classe Shape. Elle implémente la méthode area pour calculer la surface d'un cercle en fonction du rayon.
    classe Rectangle(Forme) : Cette classe définit une autre sous-classe Rectangle de la classe Shape. Elle met en œuvre la méthode area pour calculer la surface d'un rectangle en fonction de sa longueur et de sa largeur.
    Les instances des classes Cercle et Rectangle sont créées en invitant l'utilisateur à saisir les dimensions requises.
    La méthode area est appelée sur les objets Circle et Rectangle pour calculer et imprimer les surfaces du cercle et du rectangle, respectivement.

Voici ce que fait le code :

    Il définit une classe de base abstraite Shape avec une méthode area abstraite.
    Il crée deux sous-classes concrètes, Circle et Rectangle, qui héritent de Shape et implémentent la méthode area.
    Elle invite l'utilisateur à saisir les dimensions d'un cercle et d'un rectangle.
    Il calcule et imprime les aires du cercle et du rectangle à l'aide des méthodes d'aire respectives.
        \end{solution}
        

        \question
        Ecrire un programme python pour convertir les heures en jours
        \par
        \begin{solution}
            \renewcommand{\nomfichier}{q738.py}
            \pythonfile{\chemincode \nomfichier}[][\nomfichier]
            Le programme python définit une classe HoursToDaysConverter, qui est utilisée pour convertir un nombre d'heures en jours. Voici la décomposition du code :

    class HoursToDaysConverter : Cette classe représente un convertisseur permettant de convertir des heures en jours. Elle possède une méthode __init__ pour initialiser l'attribut hours et une méthode convert_to_days pour effectuer la conversion.
    Dans le bloc try, le code tente d'obtenir le nombre d'heures de l'utilisateur à l'aide de la fonction input. Il convertit l'entrée de l'utilisateur en un nombre flottant et l'affecte à la variable hours_input. Si l'utilisateur saisit une valeur non valide (qui n'est pas un nombre), une exception ValueError est levée et un message d'erreur est imprimé.
    Si l'utilisateur saisit un nombre d'heures valide, une instance de la classe HoursToDaysConverter est créée à partir de la variable hours_input et la méthode convert_to_days est appelée pour effectuer la conversion.
    Le résultat, qui est le nombre de jours, est stocké dans la variable days_result.
    Enfin, le code imprime le nombre original d'heures saisi par l'utilisateur et le nombre équivalent de jours basé sur la conversion.

Voici ce que fait le code :

    Il prend le nombre d'heures saisi par l'utilisateur et tente de le convertir en un nombre flottant.
    Si l'entrée est valide, il convertit le nombre d'heures en jours à l'aide de la classe HoursToDaysConverter.
    Il affiche ensuite le nombre d'heures original et le nombre de jours équivalent.
        \end{solution}
        

        \question
        Écrire un programme python pour rechercher des objets dans un tableau d'objets à l'aide de la méthode filter().
        \par
        \begin{solution}
            \renewcommand{\nomfichier}{q739.py}
            \pythonfile{\chemincode \nomfichier}[][\nomfichier]
            Le programme python définit une classe Personne et une liste d'objets Personne. Il utilise ensuite la fonction filter() pour filtrer les personnes de la liste qui ont plus de 25 ans. Voici la décomposition du code :

    classe Personne : Cette classe définit une personne avec des attributs de nom et d'âge.
    Le code crée une liste d'objets Personne, les personnes, avec des noms et des âges différents.
    def is_older_than_25(person) : Cette fonction définit un critère de filtrage. Elle renvoie True si l'âge de la personne est supérieur à 25 ; sinon, elle renvoie False.
    filter(is_older_than_25, people) : La fonction filter() prend la fonction is_older_than_25 et la liste des personnes, et renvoie un itérable contenant uniquement les objets Person pour lesquels la fonction is_older_than_25 renvoie True.
    liste_personnes_filtrées = liste(personnes_filtrées) : Cette ligne convertit l'itérable filtré en liste. Bien que cette conversion soit facultative, elle est effectuée ici pour faciliter la lecture en boucle des résultats filtrés à plusieurs reprises.
    Le code parcourt ensuite la liste filtrée et imprime le nom et l'âge des personnes âgées de plus de 25 ans.

Voici ce que fait le code :

    Il définit une classe Personne pour représenter les personnes avec leur nom et leur âge.
    Il crée une liste d'objets Personne.
    Il utilise la fonction filter() pour filtrer et extraire les personnes âgées de plus de 25 ans.
    Il convertit le résultat du filtrage en une liste (facultatif).
    Il affiche les noms et les âges des personnes âgées de plus de 25 ans.
        \end{solution}
        

%	
%\end{questions}





\end{document}